% Expanded finite-lattice verification reference, 29 September 2026.
% Existing data transcribed; no new numerical or VOA replay is claimed.
\section{Expanded finite-lattice reference: coordinates, ranks, and recipes}
\label{expfinite:reference}

This section brings the finite construction's scattered ingredients into
one reading order. It prints the complete nine weight-two input fields,
the short-generation recipe graph, the grade-three weight census, and
the additional Gate~4 test list. It also records what the historical
numerical calculations found before exact certificates became available.
The arithmetic tables are transcriptions of saved records, with rational
coefficients rewritten in the basis $1,s,t,st$, where $s^2=13$ and
$t^2=-3$. Their inclusion here is a documentation step, not a fresh replay
of the vertex-mode calculations.

The distinction between three kinds of assertion is used throughout.
An \emph{ambient identity} compares actual lattice vectors. A
\emph{generation recipe} establishes that such a vector is in the VOA
generated by the six currents, $\omega_B$, and $u$. A
\emph{closure argument} shows that all higher products stay in the
normal span. The exact fields alone supply none of the last two
implications. Conversely, a successful normal-space membership test
need not compute the full dimension of the ambient normal space.

\subsection{A fully specified coordinate system}
\label{expfinite:coordinates}

For direct checking the Gram matrix is
\[
G=\begin{pmatrix}
2&-1&0&0&0&0&0&0\\
-1&2&-1&0&0&0&0&0\\
0&-1&2&-1&0&-1&0&0\\
0&0&-1&2&-1&0&0&0\\
0&0&0&-1&2&0&0&0\\
0&0&-1&0&0&2&0&0\\
0&0&0&0&0&0&2&-1\\
0&0&0&0&0&0&-1&2
\end{pmatrix}.
\]
All momenta are columns in this ordered simple-root basis, even though
the tuple notation suppresses transposition. Thus
$(\alpha,\beta)=\alpha^{\mathsf T}G\beta$. Let
\[
 [\lambda;\mathcal O]
 =\prod_{k\in\mathcal O}\alpha_{k\bmod8}
       (-1-\lfloor k/8\rfloor)e^\lambda,
\qquad
\wt[\lambda;\mathcal O]
 =\tfrac12\lambda^{\mathsf T}G\lambda+
       \sum_{k\in\mathcal O}(1+\lfloor k/8\rfloor).
\]
The oscillator list is sorted and repeated entries are retained.
The empty list means the pure exponential $e^\lambda$. These are
unnormalized simple-root oscillators. In particular,
$[0;(i,i)]$ has no hidden factor of $\sqrt2$ or $1/2$.

Two practical examples fix the mode conventions:
\[
 [0;(0,8)]=\alpha_0(-1)\alpha_0(-2)\mathbf1,
 \qquad [\lambda;(6)]=\alpha_6(-1)e^\lambda.
\]
In the exact engine, a coefficient encoded by
$[[q_0,q_s],[q_t,q_{st}]]$ means
$q_0+q_ss+q_tt+q_{st}st$. The two inner lists separate the
$rational+s$ and $t+st$ parts; they are not numerical real and
imaginary coordinates. Decimal fields retained in the old JSON are
historical reconstruction diagnostics and are omitted from the exact
tables below.

For an independent implementation, all required arithmetic can be done
with rational fractions. If $x=x_0+x_1s+x_2t+x_3st$ and
$y=y_0+y_1s+y_2t+y_3st$, then
\begin{align*}
(xy)_0&=x_0y_0+13x_1y_1-3x_2y_2-39x_3y_3,\\
(xy)_1&=x_0y_1+x_1y_0-3x_2y_3-3x_3y_2,\\
(xy)_2&=x_0y_2+x_2y_0+13x_1y_3+13x_3y_1,\\
(xy)_3&=x_0y_3+x_3y_0+x_1y_2+x_2y_1.
\end{align*}
Inversion can be performed by multiplying the other three Galois
conjugates and dividing by their product with $x$. Zero testing is four
rational zero tests. Complex conjugation is the automorphism fixing $s$
and negating $t$; it must still be combined with momentum reversal and
the PCT gauge factor when applied to a state.

\subsection{Reconstructing an arbitrary mode from the printed inputs}
\label{expfinite:mode-engine}

The root cocycle in these coordinates is
\[
 \epsilon(\alpha,\beta)=(-1)^{
 \sum_i\alpha_i\beta_i+\sum_{i>j}\alpha_i\beta_jG_{ij}}.
\]
Its exponent is bilinear modulo two, so associativity of the twisted
lattice multiplication follows directly. For a root $\alpha$,
$\epsilon(\alpha,-\alpha)=-1$; this accounts for the minus sign in
both lowering-current formulas. A checker that uses the opposite root
sign and leaves all field coefficients unchanged is checking a different
coordinate convention.

For a pure exponential acting on an arbitrary oscillator state, the
following finite coefficient extraction is sufficient:
\begin{align}
Y(e^\alpha,z)\prod_{j=1}^r h_j(-n_j)e^\beta
={}&\epsilon(\alpha,\beta)z^{(\alpha,\beta)}
\prod_{j=1}^r\bigl(h_j(-n_j)-(h_j,\alpha)z^{-n_j}\bigr)
\notag\\[-2pt]
&\hspace{10mm}\cdot
\exp\!\left(\sum_{n>0}\frac{\alpha(-n)}n z^n\right)
 e^{\alpha+\beta}.
\label{expfinite:full-vertex-action}
\end{align}
To obtain $(e^\alpha)_m$, keep the coefficient of $z^{-m-1}$.
Expanding the finite product corresponds to choosing a subset of the
input oscillators to contract. If $S$ is the contracted subset, the
required degree in the creation exponential is
\[
 d_S=-m-1-(\alpha,\beta)+\sum_{j\in S}n_j.
\]
Terms with $d_S<0$ are absent. At a fixed target conformal weight there
are only finitely many partitions of $d_S$ and finitely many momenta.
This is the mathematical meaning of the exact engine's contraction-mask
loop and its creation-polynomial cache.

The first creation polynomials, useful for manually checking the
weight-two rows, are
\begin{align*}
C_0(\alpha)&=1,\\
C_1(\alpha)&=\alpha(-1),\\
C_2(\alpha)&=\tfrac12\alpha(-2)+\tfrac12\alpha(-1)^2,\\
C_3(\alpha)&=\tfrac13\alpha(-3)+\tfrac12\alpha(-2)\alpha(-1)
                  +\tfrac16\alpha(-1)^3.
\end{align*}
In simple-root coordinates the mixed coefficient of
$\alpha_i(-1)\alpha_j(-1)$ in $C_2$ is $\alpha_i\alpha_j$ for
$i\ne j$, and the diagonal coefficient is $\alpha_i^2/2$.

Heisenberg modes are even simpler. For $m>0$, $h(m)$ contracts each
occurrence of $k(-m)$ with coefficient $m(h,k)$, including its
multiplicity. The zero mode is multiplication by $(h,\beta)$, and
negative modes insert another oscillator. General oscillator fields
are obtained from normal products and derivatives:
\[
Y(h(-n)a,z)=\frac1{(n-1)!}:
   (\partial_z^{n-1}h(z))Y(a,z):,
\qquad
(Ta)_m=-m a_{m-1}.
\]
The convention $a_mb\in B_{\wt a+\wt b-m-1}$ provides a check on
every intermediate state. No matrices for all ambient grades are
required to interpret this recipe.

\subsection{Ambient census and the affine search space}
\label{expfinite:ambient-census}

The saved lattice enumeration finds $78$ norm-two vectors, $702$
norm-four vectors, and $2346$ norm-six vectors. The ambient graded
dimensions are
\[
 \dim B_0=1,\quad\dim B_1=86,\quad
 \dim B_2=1370,\quad\dim B_3=11586.
\]
The first two nontrivial counts can be checked without a VOA product:
\[
86=8+78,
\qquad
1370=(8+36)+78\cdot8+702.
\]
At weight three the rank-eight oscillator counts are
$8+64+120=192$, so
\[
11586=192+78\cdot44+702\cdot8+2346.
\]
Here $44=8+36$ is the number of oscillator states of degree two;
$120$ is the number of three-oscillator monomials of mode number one.
These are ambient counts, unrelated to the later $120$-dimensional
selected space except through inclusion.

The original highest-space search imposed the two raising zero modes
and all six positive current modes at weight two. Its saved list of
ambient affine-primary multiplicities is printed below. The entries
$(n,m)$ are representation dimensions; the Cartan eigenvalues of a
highest vector are $(n-1,m-1)$. The last column is the dimension of the
ambient joint Cartan eigenspace before imposing the kernel equations.
This table is a floating kernel search result, from
\artifact{hi_cft_20260923_c8_lattice_control.json}; it is useful for
locating the exact selected fields, but its multiplicities are not
being promoted here to a new exact certificate.

\begin{longtable}{@{}crr@{}}
\toprule Type $(n,m)$&Primary multiplicity&Cartan-block dimension\\\midrule\endfirsthead
\toprule Type $(n,m)$&Primary multiplicity&Cartan-block dimension\\\midrule\endhead
\bottomrule\endfoot
$(1,1)$ & 4 & 86\\
$(1,3)$ & 6 & 78\\
$(1,5)$ & 11 & 66\\
$(1,7)$ & 9 & 48\\
$(1,9)$ & 11 & 34\\
$(1,11)$ & 7 & 19\\
$(1,13)$ & 5 & 9\\
$(1,15)$ & 2 & 3\\
$(1,17)$ & 1 & 1\\
$(2,2)$ & 4 & 47\\
$(2,4)$ & 5 & 42\\
$(2,6)$ & 8 & 34\\
$(2,8)$ & 7 & 23\\
$(2,10)$ & 5 & 14\\
$(2,12)$ & 3 & 7\\
$(2,14)$ & 3 & 3\\
\end{longtable}
The same exploratory record gives a largest current-algebra and
adjoint residual of approximately $1.16\times10^{-14}$. This measures
the numerical realization in its original orthonormal basis, whereas
the subsequent exact current formulas use the unnormalized simple-root
basis and the transported torus gauge.

\subsection{Complete charged highest fields and retained neutral stress}
\label{expfinite:complete-fields}

The next nine tables give every nonzero coordinate of the eight
charged highest fields $a_0,\ldots,a_7$ and the retained neutral stress
$e$. A row $(\lambda,\mathcal O,c)$ means $c[\lambda;\mathcal O]$,
and the field is the sum of all its rows. They therefore specify actual
ambient lattice vectors without opening an NPZ file or a recognition
script. The second neutral field is recovered exactly as
\[
 f=\omega_B-\omega_{A_{1,1}}-\omega_{A_{1,39}}-e.
\]
Its independent coordinate table is unnecessary once the currents,
ambient stress, and $e$ are given.

The source is \artifact{hi_cft_20260923_c8_exact_fields.json}.
That file retains its original historical scope string, ``Floating
reconstruction conjectures, requiring subsequent exact verification.''
The later exact and exact-unitary certificates verify the corrected
rational entries as lattice fields. In particular, the recorded repair
changes the coefficients of the two neutral oscillator states
$[0;(0,7)]$ and $[0;(4,6)]$ from
\[
 \frac{1799-503s}{4085}
 \quad\hbox{to}\quad
 -\frac{260}{3813}+\frac{5s}{279}.
\]
The tables below use the corrected entries, not the saved numerical
recognition approximations. The repair was constrained inside the
four-dimensional commutant space; its provenance is recorded in the
source's \texttt{stress\_repair} object.

The remaining descendants of a highest field of type $(n,m)$ are
obtained by the specified current modes
\[
 (F_L)_0^i(F_R)_0^j a,
 \qquad 0\le i<n,\quad 0\le j<m.
\]
There are no divided powers in this recipe. This convention matters
when comparing coefficients with checkpoint vectors, which may instead
have been normalized to a pivot equal to one.

\subsubsection{\texorpdfstring{Field $a_{0}$: type $(1,5)$, 61 terms}{Field a0: type (1,5), 61 terms}}
The $e_1$ eigenvalue is $\frac{798}{1271}$. 
\begingroup\scriptsize\setlength{\tabcolsep}{3pt}\renewcommand{\arraystretch}{1.35}
\begin{longtable}{@{}llr@{}}
\toprule Momentum $\lambda$&Oscillators $\mathcal O$&Coefficient\\\midrule\endfirsthead
\toprule Momentum $\lambda$&Oscillators $\mathcal O$&Coefficient\\\midrule\endhead
\bottomrule\endfoot
$(-1,0,1,1,1,0,0,0)$ & $\varnothing$ & $\frac{1159+589st}{10920}$\\
$(0,0,0,0,0,0,1,1)$ & $(0)$ & $\frac{-57s-95t}{364}$\\
$(0,0,0,0,0,0,1,1)$ & $(1)$ & $\frac{-19s+19t}{182}$\\
$(0,0,0,0,0,0,1,1)$ & $(2)$ & $\frac{57t}{182}$\\
$(0,0,0,0,0,0,1,1)$ & $(3)$ & $\frac{19s+19t}{182}$\\
$(0,0,0,0,0,0,1,1)$ & $(4)$ & $\frac{57s-95t}{364}$\\
$(0,0,0,0,0,0,1,1)$ & $(6)$ & $-1$\\
$(0,0,0,0,0,0,1,1)$ & $(7)$ & $1$\\
$(0,0,0,0,1,0,0,1)$ & $\varnothing$ & $\frac{16+19s-57t}{182}$\\
$(0,0,0,0,1,0,1,0)$ & $\varnothing$ & $\frac{16-19s+57t}{182}$\\
$(0,0,0,1,0,0,0,1)$ & $\varnothing$ & $\frac{140+19s+57t}{728}$\\
$(0,0,0,1,0,0,1,0)$ & $\varnothing$ & $\frac{140-19s-57t}{728}$\\
$(0,0,0,1,1,0,0,0)$ & $(0)$ & $\frac{703+133st}{3640}$\\
$(0,0,0,1,1,0,0,0)$ & $(1)$ & $\frac{57+19st}{728}$\\
$(0,0,0,1,1,0,0,0)$ & $(2)$ & $\frac{-57-57st}{1820}$\\
$(0,0,0,1,1,0,0,0)$ & $(3)$ & $\frac{-167-57st}{1820}$\\
$(0,0,0,1,1,0,0,0)$ & $(4)$ & $\frac{11}{182}$\\
$(0,0,0,1,1,0,0,0)$ & $(6)$ & $\frac{-76-19s+95t}{728}$\\
$(0,0,0,1,1,0,0,0)$ & $(7)$ & $\frac{-76+19s-95t}{728}$\\
$(0,0,1,0,0,0,0,1)$ & $\varnothing$ & $\frac{62+57t}{273}$\\
$(0,0,1,0,0,0,1,0)$ & $\varnothing$ & $\frac{62-57t}{273}$\\
$(0,0,1,0,1,0,0,0)$ & $\varnothing$ & $\frac{-89-19st}{1365}$\\
$(0,0,1,1,0,0,0,0)$ & $(0)$ & $\frac{95-19st}{1092}$\\
$(0,0,1,1,0,0,0,0)$ & $(1)$ & $\frac{57+19st}{728}$\\
$(0,0,1,1,0,0,0,0)$ & $(2)$ & $\frac{13+19st}{728}$\\
$(0,0,1,1,0,0,0,0)$ & $(3)$ & $\frac{-29-19st}{1092}$\\
$(0,0,1,1,0,0,0,0)$ & $(4)$ & $\frac{-95-19st}{1092}$\\
$(0,0,1,1,0,0,0,0)$ & $(6)$ & $\frac{-76-19s+95t}{2184}$\\
$(0,0,1,1,0,0,0,0)$ & $(7)$ & $\frac{-76+19s-95t}{2184}$\\
$(0,0,1,1,1,0,-1,0)$ & $\varnothing$ & $\frac{-76-19s+95t}{1092}$\\
$(0,0,1,1,1,0,0,-1)$ & $\varnothing$ & $\frac{-76+19s-95t}{1092}$\\
$(0,1,0,0,0,0,0,1)$ & $\varnothing$ & $\frac{140-19s+57t}{728}$\\
$(0,1,0,0,0,0,1,0)$ & $\varnothing$ & $\frac{140+19s-57t}{728}$\\
$(0,1,0,0,1,0,0,0)$ & $\varnothing$ & $\frac{111+171st}{3640}$\\
$(0,1,0,1,0,0,0,0)$ & $\varnothing$ & $\frac{1}{28}$\\
$(0,1,1,0,0,0,0,0)$ & $(0)$ & $\frac{95-19st}{1092}$\\
$(0,1,1,0,0,0,0,0)$ & $(1)$ & $\frac{29-19st}{1092}$\\
$(0,1,1,0,0,0,0,0)$ & $(2)$ & $\frac{-13+19st}{728}$\\
$(0,1,1,0,0,0,0,0)$ & $(3)$ & $\frac{-57+19st}{728}$\\
$(0,1,1,0,0,0,0,0)$ & $(4)$ & $\frac{-95-19st}{1092}$\\
$(0,1,1,0,0,0,0,0)$ & $(6)$ & $\frac{76-19s-95t}{2184}$\\
$(0,1,1,0,0,0,0,0)$ & $(7)$ & $\frac{76+19s+95t}{2184}$\\
$(0,1,1,1,0,0,-1,0)$ & $\varnothing$ & $\frac{-76+95t}{1092}$\\
$(0,1,1,1,0,0,0,-1)$ & $\varnothing$ & $\frac{-76-95t}{1092}$\\
$(0,1,1,1,1,0,-1,-1)$ & $\varnothing$ & $\frac{19s}{1092}$\\
$(1,0,0,0,0,0,0,1)$ & $\varnothing$ & $\frac{16-19s-57t}{182}$\\
$(1,0,0,0,0,0,1,0)$ & $\varnothing$ & $\frac{16+19s+57t}{182}$\\
$(1,0,0,0,1,0,0,0)$ & $\varnothing$ & $\frac{22}{65}$\\
$(1,0,0,1,0,0,0,0)$ & $\varnothing$ & $\frac{111-171st}{3640}$\\
$(1,0,1,0,0,0,0,0)$ & $\varnothing$ & $\frac{-89+19st}{1365}$\\
$(1,1,0,0,0,0,0,0)$ & $(0)$ & $\frac{-11}{182}$\\
$(1,1,0,0,0,0,0,0)$ & $(1)$ & $\frac{167-57st}{1820}$\\
$(1,1,0,0,0,0,0,0)$ & $(2)$ & $\frac{57-57st}{1820}$\\
$(1,1,0,0,0,0,0,0)$ & $(3)$ & $\frac{-57+19st}{728}$\\
$(1,1,0,0,0,0,0,0)$ & $(4)$ & $\frac{-703+133st}{3640}$\\
$(1,1,0,0,0,0,0,0)$ & $(6)$ & $\frac{76-19s-95t}{728}$\\
$(1,1,0,0,0,0,0,0)$ & $(7)$ & $\frac{76+19s+95t}{728}$\\
$(1,1,1,0,-1,0,0,0)$ & $\varnothing$ & $\frac{1159-589st}{10920}$\\
$(1,1,1,0,0,0,-1,0)$ & $\varnothing$ & $\frac{-76+19s+95t}{1092}$\\
$(1,1,1,0,0,0,0,-1)$ & $\varnothing$ & $\frac{-76-19s-95t}{1092}$\\
$(1,1,1,1,0,0,-1,-1)$ & $\varnothing$ & $\frac{19s}{1092}$\\
\end{longtable}
\endgroup
\subsubsection{\texorpdfstring{Field $a_{1}$: type $(1,9)$, 32 terms}{Field a1: type (1,9), 32 terms}}
The $e_1$ eigenvalue is $\frac{200}{1271}$. 
\begingroup\scriptsize\setlength{\tabcolsep}{3pt}\renewcommand{\arraystretch}{1.35}
\begin{longtable}{@{}llr@{}}
\toprule Momentum $\lambda$&Oscillators $\mathcal O$&Coefficient\\\midrule\endfirsthead
\toprule Momentum $\lambda$&Oscillators $\mathcal O$&Coefficient\\\midrule\endhead
\bottomrule\endfoot
$(0,0,0,1,1,0,1,1)$ & $\varnothing$ & $\frac{1-4st}{25}$\\
$(0,0,1,1,0,0,1,1)$ & $\varnothing$ & $\frac{-3-13st}{360}$\\
$(0,0,1,1,1,0,0,1)$ & $\varnothing$ & $\frac{-87-153s+581t+223st}{3600}$\\
$(0,0,1,1,1,0,1,0)$ & $\varnothing$ & $\frac{87-153s+581t-223st}{3600}$\\
$(0,0,1,2,1,0,0,0)$ & $\varnothing$ & $\frac{69s+287t}{3600}$\\
$(0,1,1,0,0,0,1,1)$ & $\varnothing$ & $\frac{81+st}{360}$\\
$(0,1,1,1,0,0,0,1)$ & $\varnothing$ & $\frac{42-18s+36t+7st}{360}$\\
$(0,1,1,1,0,0,1,0)$ & $\varnothing$ & $\frac{-42-18s+36t-7st}{360}$\\
$(0,1,1,1,1,0,0,0)$ & $(0)$ & $\frac{-91s-143t}{3600}$\\
$(0,1,1,1,1,0,0,0)$ & $(1)$ & $\frac{19s-363t}{3600}$\\
$(0,1,1,1,1,0,0,0)$ & $(2)$ & $\frac{33s-91t}{1200}$\\
$(0,1,1,1,1,0,0,0)$ & $(3)$ & $\frac{29s+117t}{3600}$\\
$(0,1,1,1,1,0,0,0)$ & $(4)$ & $\frac{-11s+22t}{360}$\\
$(0,1,1,1,1,0,0,0)$ & $(6)$ & $\frac{169+27s+221t-51st}{3600}$\\
$(0,1,1,1,1,0,0,0)$ & $(7)$ & $\frac{-169+27s+221t+51st}{3600}$\\
$(0,1,2,1,0,0,0,0)$ & $\varnothing$ & $\frac{-s+2t}{24}$\\
$(1,0,1,1,1,0,0,0)$ & $\varnothing$ & $\frac{-9s-107t}{900}$\\
$(1,1,0,0,0,0,1,1)$ & $\varnothing$ & $1$\\
$(1,1,0,1,1,0,0,0)$ & $\varnothing$ & $\frac{13s-26t}{200}$\\
$(1,1,1,0,0,0,0,1)$ & $\varnothing$ & $\frac{279-57s+59t-st}{720}$\\
$(1,1,1,0,0,0,1,0)$ & $\varnothing$ & $\frac{-279-57s+59t+st}{720}$\\
$(1,1,1,0,1,0,0,0)$ & $\varnothing$ & $\frac{51s+23t}{900}$\\
$(1,1,1,1,0,0,0,0)$ & $(0)$ & $\frac{11s-22t}{360}$\\
$(1,1,1,1,0,0,0,0)$ & $(1)$ & $\frac{11s+13t}{720}$\\
$(1,1,1,1,0,0,0,0)$ & $(2)$ & $\frac{-9s+13t}{240}$\\
$(1,1,1,1,0,0,0,0)$ & $(3)$ & $\frac{-7s+t}{144}$\\
$(1,1,1,1,0,0,0,0)$ & $(4)$ & $\frac{-13s-39t}{720}$\\
$(1,1,1,1,0,0,0,0)$ & $(6)$ & $\frac{65+21s+13t-5st}{720}$\\
$(1,1,1,1,0,0,0,0)$ & $(7)$ & $\frac{-65+21s+13t+5st}{720}$\\
$(1,1,1,1,1,0,-1,0)$ & $\varnothing$ & $\frac{-78-39s+78t-13st}{1800}$\\
$(1,1,1,1,1,0,0,-1)$ & $\varnothing$ & $\frac{78-39s+78t+13st}{1800}$\\
$(1,2,1,0,0,0,0,0)$ & $\varnothing$ & $\frac{-27s-31t}{720}$\\
\end{longtable}
\endgroup
\subsubsection{\texorpdfstring{Field $a_{2}$: type $(1,9)$, 32 terms}{Field a2: type (1,9), 32 terms}}
The $e_1$ eigenvalue is $\frac{1512}{1271}$. 
\begingroup\scriptsize\setlength{\tabcolsep}{3pt}\renewcommand{\arraystretch}{1.35}
\begin{longtable}{@{}llr@{}}
\toprule Momentum $\lambda$&Oscillators $\mathcal O$&Coefficient\\\midrule\endfirsthead
\toprule Momentum $\lambda$&Oscillators $\mathcal O$&Coefficient\\\midrule\endhead
\bottomrule\endfoot
$(0,0,0,1,1,0,1,1)$ & $\varnothing$ & $\frac{-7+3st}{20}$\\
$(0,0,1,1,0,0,1,1)$ & $\varnothing$ & $\frac{15+5st}{48}$\\
$(0,0,1,1,1,0,0,1)$ & $\varnothing$ & $\frac{159-43s+11t-11st}{480}$\\
$(0,0,1,1,1,0,1,0)$ & $\varnothing$ & $\frac{-159-43s+11t+11st}{480}$\\
$(0,0,1,2,1,0,0,0)$ & $\varnothing$ & $\frac{11s+53t}{480}$\\
$(0,1,1,0,0,0,1,1)$ & $\varnothing$ & $\frac{6+st}{12}$\\
$(0,1,1,1,0,0,0,1)$ & $\varnothing$ & $\frac{9-2s-6t-st}{96}$\\
$(0,1,1,1,0,0,1,0)$ & $\varnothing$ & $\frac{-9-2s-6t+st}{96}$\\
$(0,1,1,1,1,0,0,0)$ & $(0)$ & $\frac{s-2t}{24}$\\
$(0,1,1,1,1,0,0,0)$ & $(1)$ & $\frac{3s-11t}{96}$\\
$(0,1,1,1,1,0,0,0)$ & $(2)$ & $\frac{s-2t}{40}$\\
$(0,1,1,1,1,0,0,0)$ & $(3)$ & $\frac{7s+11t}{240}$\\
$(0,1,1,1,1,0,0,0)$ & $(4)$ & $\frac{s+3t}{96}$\\
$(0,1,1,1,1,0,0,0)$ & $(6)$ & $\frac{-38-33s+41t+2st}{480}$\\
$(0,1,1,1,1,0,0,0)$ & $(7)$ & $\frac{38-33s+41t-2st}{480}$\\
$(0,1,2,1,0,0,0,0)$ & $\varnothing$ & $\frac{s+3t}{96}$\\
$(1,0,1,1,1,0,0,0)$ & $\varnothing$ & $\frac{s-t}{24}$\\
$(1,1,0,0,0,0,1,1)$ & $\varnothing$ & $1$\\
$(1,1,0,1,1,0,0,0)$ & $\varnothing$ & $\frac{-s-3t}{160}$\\
$(1,1,1,0,0,0,0,1)$ & $\varnothing$ & $\frac{6+s-4t-st}{24}$\\
$(1,1,1,0,0,0,1,0)$ & $\varnothing$ & $\frac{-6+s-4t+st}{24}$\\
$(1,1,1,0,1,0,0,0)$ & $\varnothing$ & $\frac{-s+2t}{30}$\\
$(1,1,1,1,0,0,0,0)$ & $(0)$ & $\frac{-s-3t}{96}$\\
$(1,1,1,1,0,0,0,0)$ & $(1)$ & $\frac{-s-7t}{96}$\\
$(1,1,1,1,0,0,0,0)$ & $(2)$ & $\frac{s-t}{32}$\\
$(1,1,1,1,0,0,0,0)$ & $(3)$ & $\frac{3s-t}{48}$\\
$(1,1,1,1,0,0,0,0)$ & $(4)$ & $\frac{5s-5t}{96}$\\
$(1,1,1,1,0,0,0,0)$ & $(6)$ & $\frac{5-6s+10t-st}{96}$\\
$(1,1,1,1,0,0,0,0)$ & $(7)$ & $\frac{-5-6s+10t+st}{96}$\\
$(1,1,1,1,1,0,-1,0)$ & $\varnothing$ & $\frac{-63-3s-9t+7st}{480}$\\
$(1,1,1,1,1,0,0,-1)$ & $\varnothing$ & $\frac{63-3s-9t-7st}{480}$\\
$(1,2,1,0,0,0,0,0)$ & $\varnothing$ & $\frac{s+2t}{24}$\\
\end{longtable}
\endgroup
\subsubsection{\texorpdfstring{Field $a_{3}$: type $(1,13)$, 9 terms}{Field a3: type (1,13), 9 terms}}
The $e_1$ eigenvalue is $\frac{420}{1271}$. 
\begingroup\scriptsize\setlength{\tabcolsep}{3pt}\renewcommand{\arraystretch}{1.35}
\begin{longtable}{@{}llr@{}}
\toprule Momentum $\lambda$&Oscillators $\mathcal O$&Coefficient\\\midrule\endfirsthead
\toprule Momentum $\lambda$&Oscillators $\mathcal O$&Coefficient\\\midrule\endhead
\bottomrule\endfoot
$(0,1,1,1,1,0,1,1)$ & $\varnothing$ & $\frac{6+2s-4t+st}{14}$\\
$(0,1,2,2,1,0,0,0)$ & $\varnothing$ & $\frac{-26+9s-13t-6st}{168}$\\
$(1,1,1,1,0,0,1,1)$ & $\varnothing$ & $\frac{-6+2s+4t+st}{14}$\\
$(1,1,1,1,1,0,0,1)$ & $\varnothing$ & $\frac{1+4t}{7}$\\
$(1,1,1,1,1,0,1,0)$ & $\varnothing$ & $1$\\
$(1,1,1,2,1,0,0,0)$ & $\varnothing$ & $\frac{-2-3s-t+2st}{56}$\\
$(1,1,2,1,1,0,0,0)$ & $\varnothing$ & $\frac{-4-2t}{21}$\\
$(1,2,1,1,1,0,0,0)$ & $\varnothing$ & $\frac{-2+3s-t-2st}{56}$\\
$(1,2,2,1,0,0,0,0)$ & $\varnothing$ & $\frac{-26-9s-13t+6st}{168}$\\
\end{longtable}
\endgroup
\subsubsection{\texorpdfstring{Field $a_{4}$: type $(2,2)$, 42 terms}{Field a4: type (2,2), 42 terms}}
The $e_1$ eigenvalue is $\frac{192}{1271}$. 
\begingroup\scriptsize\setlength{\tabcolsep}{3pt}\renewcommand{\arraystretch}{1.35}
\begin{longtable}{@{}llr@{}}
\toprule Momentum $\lambda$&Oscillators $\mathcal O$&Coefficient\\\midrule\endfirsthead
\toprule Momentum $\lambda$&Oscillators $\mathcal O$&Coefficient\\\midrule\endhead
\bottomrule\endfoot
$(-1,-2,-2,-1,0,-1,1,1)$ & $\varnothing$ & $\frac{217+32st}{295}$\\
$(-1,-2,-1,0,0,-1,0,0)$ & $\varnothing$ & $\frac{-29s-42t}{885}$\\
$(-1,-1,-2,-1,-1,-1,1,1)$ & $\varnothing$ & $\frac{32+2st}{59}$\\
$(-1,-1,-2,-1,0,-1,0,1)$ & $\varnothing$ & $\frac{114-15s+80t+44st}{1180}$\\
$(-1,-1,-2,-1,0,-1,1,0)$ & $\varnothing$ & $\frac{-114-15s+80t-44st}{1180}$\\
$(-1,-1,-2,0,0,-1,0,0)$ & $\varnothing$ & $\frac{-29s-42t}{1180}$\\
$(-1,-1,-1,-1,-1,-1,0,1)$ & $\varnothing$ & $\frac{192-15s+80t+12st}{708}$\\
$(-1,-1,-1,-1,-1,-1,1,0)$ & $\varnothing$ & $\frac{-192-15s+80t-12st}{708}$\\
$(-1,-1,-1,-1,0,-1,0,0)$ & $(1)$ & $\frac{-29s-42t}{885}$\\
$(-1,-1,-1,-1,0,-1,0,0)$ & $(2)$ & $\frac{-29s-42t}{1180}$\\
$(-1,-1,-1,-1,0,-1,0,0)$ & $(3)$ & $\frac{22s-19t}{885}$\\
$(-1,-1,-1,-1,0,-1,0,0)$ & $(4)$ & $\frac{35s+10t}{708}$\\
$(-1,-1,-1,-1,0,-1,0,0)$ & $(6)$ & $\frac{-103-15s+80t+12st}{1770}$\\
$(-1,-1,-1,-1,0,-1,0,0)$ & $(7)$ & $\frac{103-15s+80t-12st}{1770}$\\
$(-1,-1,-1,0,-1,-1,0,0)$ & $\varnothing$ & $\frac{35s+10t}{708}$\\
$(-1,-1,-1,0,0,-1,-1,0)$ & $\varnothing$ & $\frac{-103-15s+80t+12st}{1770}$\\
$(-1,-1,-1,0,0,-1,0,-1)$ & $\varnothing$ & $\frac{103-15s+80t-12st}{1770}$\\
$(-1,0,-1,-1,-1,-1,0,0)$ & $\varnothing$ & $\frac{-29s-42t}{708}$\\
$(0,-1,-2,-2,-1,-1,1,1)$ & $\varnothing$ & $1$\\
$(0,-1,-2,-1,-1,-1,0,1)$ & $\varnothing$ & $\frac{-54+3s-16t+4st}{236}$\\
$(0,-1,-2,-1,-1,-1,1,0)$ & $\varnothing$ & $\frac{54+3s-16t-4st}{236}$\\
$(0,-1,-2,-1,0,-1,0,0)$ & $(0)$ & $\frac{87s+126t}{1180}$\\
$(0,-1,-2,-1,0,-1,0,0)$ & $(1)$ & $\frac{13s+29t}{295}$\\
$(0,-1,-2,-1,0,-1,0,0)$ & $(2)$ & $\frac{-9s+48t}{590}$\\
$(0,-1,-2,-1,0,-1,0,0)$ & $(3)$ & $\frac{-19s+3t}{295}$\\
$(0,-1,-2,-1,0,-1,0,0)$ & $(4)$ & $\frac{-21s-6t}{236}$\\
$(0,-1,-2,-1,0,-1,0,0)$ & $(6)$ & $\frac{64+15s-80t+4st}{590}$\\
$(0,-1,-2,-1,0,-1,0,0)$ & $(7)$ & $\frac{-64+15s-80t-4st}{590}$\\
$(0,-1,-1,-1,-1,-1,0,0)$ & $(0)$ & $\frac{29s+42t}{708}$\\
$(0,-1,-1,-1,-1,-1,0,0)$ & $(1)$ & $\frac{2s+9t}{177}$\\
$(0,-1,-1,-1,-1,-1,0,0)$ & $(2)$ & $\frac{-7s-2t}{236}$\\
$(0,-1,-1,-1,-1,-1,0,0)$ & $(3)$ & $\frac{-7s-2t}{177}$\\
$(0,-1,-1,-1,-1,-1,0,0)$ & $(6)$ & $\frac{5+3s-16t+4st}{354}$\\
$(0,-1,-1,-1,-1,-1,0,0)$ & $(7)$ & $\frac{-5+3s-16t-4st}{354}$\\
$(0,-1,-1,-1,0,-1,-1,0)$ & $\varnothing$ & $\frac{-64-15s+80t-4st}{1770}$\\
$(0,-1,-1,-1,0,-1,0,-1)$ & $\varnothing$ & $\frac{64-15s+80t+4st}{1770}$\\
$(0,-1,-1,0,0,-1,-1,-1)$ & $\varnothing$ & $\frac{-39+16st}{1770}$\\
$(0,0,-2,-1,-1,-1,0,0)$ & $\varnothing$ & $\frac{-7s-2t}{236}$\\
$(0,0,-1,-2,-1,-1,0,0)$ & $\varnothing$ & $\frac{7s+2t}{177}$\\
$(0,0,-1,-1,-1,-1,-1,0)$ & $\varnothing$ & $\frac{-5-3s+16t-4st}{354}$\\
$(0,0,-1,-1,-1,-1,0,-1)$ & $\varnothing$ & $\frac{5-3s+16t+4st}{354}$\\
$(0,0,-1,-1,0,-1,-1,-1)$ & $\varnothing$ & $\frac{-39+16st}{1770}$\\
\end{longtable}
\endgroup
\subsubsection{\texorpdfstring{Field $a_{5}$: type $(2,6)$, 32 terms}{Field a5: type (2,6), 32 terms}}
The $e_1$ eigenvalue is $\frac{1830}{1271}$. 
\begingroup\scriptsize\setlength{\tabcolsep}{3pt}\renewcommand{\arraystretch}{1.35}
\begin{longtable}{@{}llr@{}}
\toprule Momentum $\lambda$&Oscillators $\mathcal O$&Coefficient\\\midrule\endfirsthead
\toprule Momentum $\lambda$&Oscillators $\mathcal O$&Coefficient\\\midrule\endhead
\bottomrule\endfoot
$(-1,-1,-1,-1,0,-1,1,1)$ & $\varnothing$ & $\frac{-55+15st}{126}$\\
$(-1,-1,-1,0,0,-1,0,1)$ & $\varnothing$ & $\frac{55-20s+40t-15st}{252}$\\
$(-1,-1,-1,0,0,-1,1,0)$ & $\varnothing$ & $\frac{-55-20s+40t+15st}{252}$\\
$(-1,-1,-1,0,1,-1,0,0)$ & $\varnothing$ & $\frac{-s+17t}{126}$\\
$(-1,-1,0,0,0,-1,0,0)$ & $\varnothing$ & $\frac{5s-5t}{252}$\\
$(-1,0,-1,0,0,-1,0,0)$ & $\varnothing$ & $\frac{5s-5t}{252}$\\
$(0,-1,-2,-1,0,-1,1,1)$ & $\varnothing$ & $1$\\
$(0,-1,-1,-1,-1,-1,1,1)$ & $\varnothing$ & $\frac{55+15st}{126}$\\
$(0,-1,-1,-1,0,-1,0,1)$ & $\varnothing$ & $\frac{4+10t}{63}$\\
$(0,-1,-1,-1,0,-1,1,0)$ & $\varnothing$ & $\frac{-4+10t}{63}$\\
$(0,-1,-1,0,0,-1,0,0)$ & $(0)$ & $\frac{-7s+31t}{252}$\\
$(0,-1,-1,0,0,-1,0,0)$ & $(1)$ & $\frac{-13s+13t}{252}$\\
$(0,-1,-1,0,0,-1,0,0)$ & $(2)$ & $\frac{-s}{21}$\\
$(0,-1,-1,0,0,-1,0,0)$ & $(3)$ & $\frac{-s-3t}{42}$\\
$(0,-1,-1,0,0,-1,0,0)$ & $(4)$ & $\frac{-s-13t}{126}$\\
$(0,-1,-1,0,0,-1,0,0)$ & $(6)$ & $\frac{13+20s-5st}{252}$\\
$(0,-1,-1,0,0,-1,0,0)$ & $(7)$ & $\frac{-13+20s+5st}{252}$\\
$(0,0,-1,-1,-1,-1,0,1)$ & $\varnothing$ & $\frac{55+20s+40t+15st}{252}$\\
$(0,0,-1,-1,-1,-1,1,0)$ & $\varnothing$ & $\frac{-55+20s+40t-15st}{252}$\\
$(0,0,-1,-1,0,-1,0,0)$ & $(0)$ & $\frac{-s+13t}{126}$\\
$(0,0,-1,-1,0,-1,0,0)$ & $(1)$ & $\frac{-s+3t}{42}$\\
$(0,0,-1,-1,0,-1,0,0)$ & $(2)$ & $\frac{-s}{21}$\\
$(0,0,-1,-1,0,-1,0,0)$ & $(3)$ & $\frac{-13s-13t}{252}$\\
$(0,0,-1,-1,0,-1,0,0)$ & $(4)$ & $\frac{-7s-31t}{252}$\\
$(0,0,-1,-1,0,-1,0,0)$ & $(6)$ & $\frac{-13+20s-5st}{252}$\\
$(0,0,-1,-1,0,-1,0,0)$ & $(7)$ & $\frac{13+20s+5st}{252}$\\
$(0,0,-1,0,-1,-1,0,0)$ & $\varnothing$ & $\frac{-5s-5t}{252}$\\
$(0,0,-1,0,0,-1,-1,0)$ & $\varnothing$ & $\frac{-13}{126}$\\
$(0,0,-1,0,0,-1,0,-1)$ & $\varnothing$ & $\frac{13}{126}$\\
$(0,0,0,-1,-1,-1,0,0)$ & $\varnothing$ & $\frac{5s+5t}{252}$\\
$(0,0,0,0,0,-1,-1,-1)$ & $\varnothing$ & $\frac{-13}{126}$\\
$(1,0,-1,-1,-1,-1,0,0)$ & $\varnothing$ & $\frac{s+17t}{126}$\\
\end{longtable}
\endgroup
\subsubsection{\texorpdfstring{Field $a_{6}$: type $(2,8)$, 22 terms}{Field a6: type (2,8), 22 terms}}
The $e_1$ eigenvalue is $\frac{752}{1271}$. 
\begingroup\scriptsize\setlength{\tabcolsep}{3pt}\renewcommand{\arraystretch}{1.35}
\begin{longtable}{@{}llr@{}}
\toprule Momentum $\lambda$&Oscillators $\mathcal O$&Coefficient\\\midrule\endfirsthead
\toprule Momentum $\lambda$&Oscillators $\mathcal O$&Coefficient\\\midrule\endhead
\bottomrule\endfoot
$(-1,-1,-1,0,0,-1,1,1)$ & $\varnothing$ & $\frac{307+297st}{1880}$\\
$(-1,0,0,0,0,-1,0,0)$ & $\varnothing$ & $\frac{-267s-541t}{5640}$\\
$(0,-1,-1,-1,0,-1,1,1)$ & $\varnothing$ & $\frac{567+77st}{940}$\\
$(0,-1,-1,0,0,-1,0,1)$ & $\varnothing$ & $\frac{-827-385s-935t+143st}{3760}$\\
$(0,-1,-1,0,0,-1,1,0)$ & $\varnothing$ & $\frac{827-385s-935t-143st}{3760}$\\
$(0,-1,-1,0,1,-1,0,0)$ & $\varnothing$ & $\frac{103s+89t}{940}$\\
$(0,-1,0,0,0,-1,0,0)$ & $\varnothing$ & $\frac{-9s+5t}{282}$\\
$(0,0,-1,-1,-1,-1,1,1)$ & $\varnothing$ & $1$\\
$(0,0,-1,-1,0,-1,0,1)$ & $\varnothing$ & $\frac{-373-253s-319t+77st}{1880}$\\
$(0,0,-1,-1,0,-1,1,0)$ & $\varnothing$ & $\frac{373-253s-319t-77st}{1880}$\\
$(0,0,-1,0,0,-1,0,0)$ & $(0)$ & $\frac{267s+541t}{5640}$\\
$(0,0,-1,0,0,-1,0,0)$ & $(1)$ & $\frac{9s-5t}{282}$\\
$(0,0,-1,0,0,-1,0,0)$ & $(2)$ & $\frac{33s-81t}{940}$\\
$(0,0,-1,0,0,-1,0,0)$ & $(3)$ & $\frac{9s-193t}{2820}$\\
$(0,0,-1,0,0,-1,0,0)$ & $(4)$ & $\frac{-15s-23t}{282}$\\
$(0,0,-1,0,0,-1,0,0)$ & $(6)$ & $\frac{-81-363s+891t-11st}{11280}$\\
$(0,0,-1,0,0,-1,0,0)$ & $(7)$ & $\frac{81-363s+891t+11st}{11280}$\\
$(0,0,0,-1,0,-1,0,0)$ & $\varnothing$ & $\frac{-9s+193t}{2820}$\\
$(0,0,0,0,-1,-1,0,0)$ & $\varnothing$ & $\frac{15s+23t}{282}$\\
$(0,0,0,0,0,-1,-1,0)$ & $\varnothing$ & $\frac{81+363s-891t+11st}{11280}$\\
$(0,0,0,0,0,-1,0,-1)$ & $\varnothing$ & $\frac{-81+363s-891t-11st}{11280}$\\
$(1,0,-1,-1,0,-1,0,0)$ & $\varnothing$ & $\frac{59s+197t}{940}$\\
\end{longtable}
\endgroup
\subsubsection{\texorpdfstring{Field $a_{7}$: type $(2,14)$, 3 terms}{Field a7: type (2,14), 3 terms}}
The $e_1$ eigenvalue is $\frac{180}{1271}$. This is the original defining field $u$.
\begingroup\scriptsize\setlength{\tabcolsep}{3pt}\renewcommand{\arraystretch}{1.35}
\begin{longtable}{@{}llr@{}}
\toprule Momentum $\lambda$&Oscillators $\mathcal O$&Coefficient\\\midrule\endfirsthead
\toprule Momentum $\lambda$&Oscillators $\mathcal O$&Coefficient\\\midrule\endhead
\bottomrule\endfoot
$(0,0,0,0,0,-1,1,1)$ & $\varnothing$ & $1$\\
$(0,0,0,1,1,-1,0,0)$ & $\varnothing$ & $\frac{s+3t}{20}$\\
$(1,1,0,0,0,-1,0,0)$ & $\varnothing$ & $\frac{s-3t}{20}$\\
\end{longtable}
\endgroup
\subsubsection{\texorpdfstring{Field $e$: type $(1,1)$, 70 terms}{Field e: type (1,1), 70 terms}}
The $e_1$ eigenvalue is $2$. 
\begingroup\scriptsize\setlength{\tabcolsep}{3pt}\renewcommand{\arraystretch}{1.35}
\begin{longtable}{@{}llr@{}}
\toprule Momentum $\lambda$&Oscillators $\mathcal O$&Coefficient\\\midrule\endfirsthead
\toprule Momentum $\lambda$&Oscillators $\mathcal O$&Coefficient\\\midrule\endhead
\bottomrule\endfoot
$(-1,-1,0,0,0,0,1,1)$ & $\varnothing$ & $\frac{10s}{93}$\\
$(-1,-1,0,1,1,0,0,0)$ & $\varnothing$ & $\frac{-10}{93}$\\
$(-1,0,0,0,0,0,0,1)$ & $\varnothing$ & $\frac{260-205s}{3813}$\\
$(-1,0,0,0,0,0,1,0)$ & $\varnothing$ & $\frac{260+205s}{3813}$\\
$(-1,0,0,0,1,0,0,0)$ & $\varnothing$ & $\frac{5}{41}$\\
$(-1,0,0,1,0,0,0,0)$ & $\varnothing$ & $\frac{55}{3813}$\\
$(-1,0,1,0,0,0,0,0)$ & $\varnothing$ & $\frac{55}{3813}$\\
$(0,-1,-1,0,0,0,1,1)$ & $\varnothing$ & $\frac{6s}{31}$\\
$(0,-1,0,0,0,0,0,1)$ & $\varnothing$ & $\frac{416-164s}{3813}$\\
$(0,-1,0,0,0,0,1,0)$ & $\varnothing$ & $\frac{416+164s}{3813}$\\
$(0,-1,0,0,1,0,0,0)$ & $\varnothing$ & $\frac{88}{3813}$\\
$(0,-1,0,1,0,0,0,0)$ & $\varnothing$ & $\frac{166}{1271}$\\
$(0,0,-1,-1,0,0,1,1)$ & $\varnothing$ & $\frac{6s}{31}$\\
$(0,0,-1,0,0,0,0,1)$ & $\varnothing$ & $\frac{156}{1271}$\\
$(0,0,-1,0,0,0,1,0)$ & $\varnothing$ & $\frac{156}{1271}$\\
$(0,0,-1,0,1,0,0,0)$ & $\varnothing$ & $\frac{33}{1271}$\\
$(0,0,0,-1,-1,0,1,1)$ & $\varnothing$ & $\frac{10s}{93}$\\
$(0,0,0,-1,0,0,0,1)$ & $\varnothing$ & $\frac{416+164s}{3813}$\\
$(0,0,0,-1,0,0,1,0)$ & $\varnothing$ & $\frac{416-164s}{3813}$\\
$(0,0,0,0,-1,0,0,1)$ & $\varnothing$ & $\frac{260+205s}{3813}$\\
$(0,0,0,0,-1,0,1,0)$ & $\varnothing$ & $\frac{260-205s}{3813}$\\
$(0,0,0,0,0,0,0,0)$ & $(0,0)$ & $\frac{890}{11439}$\\
$(0,0,0,0,0,0,0,0)$ & $(0,1)$ & $\frac{980}{11439}$\\
$(0,0,0,0,0,0,0,0)$ & $(0,2)$ & $\frac{20}{1271}$\\
$(0,0,0,0,0,0,0,0)$ & $(0,3)$ & $\frac{-455}{11439}$\\
$(0,0,0,0,0,0,0,0)$ & $(0,4)$ & $\frac{-925}{11439}$\\
$(0,0,0,0,0,0,0,0)$ & $(0,6)$ & $\frac{-780-205s}{11439}$\\
$(0,0,0,0,0,0,0,0)$ & $(0,7)$ & $\frac{-780+205s}{11439}$\\
$(0,0,0,0,0,0,0,0)$ & $(1,1)$ & $\frac{980}{11439}$\\
$(0,0,0,0,0,0,0,0)$ & $(1,2)$ & $\frac{73}{1271}$\\
$(0,0,0,0,0,0,0,0)$ & $(1,3)$ & $\frac{-646}{11439}$\\
$(0,0,0,0,0,0,0,0)$ & $(1,4)$ & $\frac{-455}{11439}$\\
$(0,0,0,0,0,0,0,0)$ & $(1,6)$ & $\frac{-1248-164s}{11439}$\\
$(0,0,0,0,0,0,0,0)$ & $(1,7)$ & $\frac{-1248+164s}{11439}$\\
$(0,0,0,0,0,0,0,0)$ & $(2,2)$ & $\frac{63}{1271}$\\
$(0,0,0,0,0,0,0,0)$ & $(2,3)$ & $\frac{73}{1271}$\\
$(0,0,0,0,0,0,0,0)$ & $(2,4)$ & $\frac{20}{1271}$\\
$(0,0,0,0,0,0,0,0)$ & $(2,6)$ & $\frac{-156}{1271}$\\
$(0,0,0,0,0,0,0,0)$ & $(2,7)$ & $\frac{-156}{1271}$\\
$(0,0,0,0,0,0,0,0)$ & $(3,3)$ & $\frac{980}{11439}$\\
$(0,0,0,0,0,0,0,0)$ & $(3,4)$ & $\frac{980}{11439}$\\
$(0,0,0,0,0,0,0,0)$ & $(3,6)$ & $\frac{-1248+164s}{11439}$\\
$(0,0,0,0,0,0,0,0)$ & $(3,7)$ & $\frac{-1248-164s}{11439}$\\
$(0,0,0,0,0,0,0,0)$ & $(4,4)$ & $\frac{890}{11439}$\\
$(0,0,0,0,0,0,0,0)$ & $(4,6)$ & $\frac{-780+205s}{11439}$\\
$(0,0,0,0,0,0,0,0)$ & $(4,7)$ & $\frac{-780-205s}{11439}$\\
$(0,0,0,0,0,0,0,0)$ & $(6,6)$ & $\frac{1820}{11439}$\\
$(0,0,0,0,0,0,0,0)$ & $(6,7)$ & $\frac{1820}{11439}$\\
$(0,0,0,0,0,0,0,0)$ & $(7,7)$ & $\frac{1820}{11439}$\\
$(0,0,0,0,1,0,-1,0)$ & $\varnothing$ & $\frac{104-82s}{3813}$\\
$(0,0,0,0,1,0,0,-1)$ & $\varnothing$ & $\frac{104+82s}{3813}$\\
$(0,0,0,1,0,0,-1,0)$ & $\varnothing$ & $\frac{104-41s}{3813}$\\
$(0,0,0,1,0,0,0,-1)$ & $\varnothing$ & $\frac{104+41s}{3813}$\\
$(0,0,0,1,1,0,-1,-1)$ & $\varnothing$ & $\frac{s}{93}$\\
$(0,0,1,0,-1,0,0,0)$ & $\varnothing$ & $\frac{55}{3813}$\\
$(0,0,1,0,0,0,-1,0)$ & $\varnothing$ & $\frac{104}{3813}$\\
$(0,0,1,0,0,0,0,-1)$ & $\varnothing$ & $\frac{104}{3813}$\\
$(0,0,1,1,0,0,-1,-1)$ & $\varnothing$ & $\frac{s}{93}$\\
$(0,1,0,-1,0,0,0,0)$ & $\varnothing$ & $\frac{166}{1271}$\\
$(0,1,0,0,-1,0,0,0)$ & $\varnothing$ & $\frac{55}{3813}$\\
$(0,1,0,0,0,0,-1,0)$ & $\varnothing$ & $\frac{104+41s}{3813}$\\
$(0,1,0,0,0,0,0,-1)$ & $\varnothing$ & $\frac{104-41s}{3813}$\\
$(0,1,1,0,0,0,-1,-1)$ & $\varnothing$ & $\frac{s}{93}$\\
$(1,0,-1,0,0,0,0,0)$ & $\varnothing$ & $\frac{33}{1271}$\\
$(1,0,0,-1,0,0,0,0)$ & $\varnothing$ & $\frac{88}{3813}$\\
$(1,0,0,0,-1,0,0,0)$ & $\varnothing$ & $\frac{5}{41}$\\
$(1,0,0,0,0,0,-1,0)$ & $\varnothing$ & $\frac{104+82s}{3813}$\\
$(1,0,0,0,0,0,0,-1)$ & $\varnothing$ & $\frac{104-82s}{3813}$\\
$(1,1,0,-1,-1,0,0,0)$ & $\varnothing$ & $\frac{-10}{93}$\\
$(1,1,0,0,0,0,-1,-1)$ & $\varnothing$ & $\frac{s}{93}$\\
\end{longtable}
\endgroup

\subsection{Exactly what the low-weight checker asserts}
\label{expfinite:low-checker}

The implementation \artifact{hi_cft_20260923_c8_exact_check.py}
reconstructs the printed fields in the exact coefficient field and
checks the following equations. With
$T_K=\omega_B-\omega_{A_{1,1}}-\omega_{A_{1,39}}$ and $f=T_K-e$,
it tests
\[
 e_1e=2e,\quad e_2e=0,\quad
 e_3e=\frac{1200}{1271}\mathbf1,
 \qquad e_nf=0\quad(n=1,2,3),\quad
 f_3f=\frac{35}{31}\mathbf1.
\]
For a highest input $a_i$ of type $(n,m)$ it tests
\begin{gather*}
(E_L)_0a_i=(E_R)_0a_i=0,\qquad
x_1a_i=0\quad(x\in K),\\
e_1a_i=h_ea_i,\quad e_2a_i=0,\quad
(\omega_B)_2a_i=0,\quad (\omega_B)_1a_i=2a_i.
\end{gather*}
The lowering strings terminate at the stated dimensions. Exact sparse
elimination finds $22$ affine normal/translation descendants, $24$
vectors after adjoining $e,f$, and $120$ after the eight multiplets.
The transported PCT fixes $e,f$ and sends every charged highest field
into this same $120$-dimensional span.

When the closure option is used, every one of the eight left highest
fields is paired with all $96$ charged descendants. For each pair
$(a_i,v)$ the checker tests both
\[
 (a_i)_1v\in S_2,\qquad (a_i)_2v\in K.
\]
Thus the eight rows labelled ``96 exact pairs'' mean $768$ pairs and
two membership tests per pair. The boolean called
\texttt{exact\_OPE\_closure} in this record has this low-weight scope;
it does not assert the later Gate~4, Gate~5, or all-weight closure.
The current, neutral-stress and descendant consequences use covariance
and the separately saved lower-closure inputs.

Three useful scalar checks on the selected stress eigenvalues are
\begin{align*}
\operatorname{Tr}_{S_2}(e_1)&=\frac{67200}{1271},\\
\operatorname{Tr}_{S_2}\bigl(e_1(H_L)_0^2/2\bigr)
   &=\frac{19900}{1271},\\
\operatorname{Tr}_{S_2}\bigl(e_1(H_R)_0^2/2\bigr)
   &=\frac{792900}{1271}.
\end{align*}
They appear in \artifact{hi_cft_20260923_c8_trace_certificate.json}.
For a current multiplet of dimensions $(n,m)$ and stress eigenvalue
$h$, its contributions to these three sums are respectively
\[
 nmh,\qquad \frac{mn(n^2-1)}6h,
 \qquad\frac{nm(m^2-1)}6h.
\]
The neutral field $e$ contributes $2$ only to the first sum; the affine
vacuum descendants and $f$ contribute zero. These formulas make all
three traces a short rational calculation from the eight eigenvalue
rows. They are necessary consistency checks and do not select the
fields uniquely or establish any OPE closure.

\subsection{The complete accepted short-generation graph}
\label{expfinite:generation-graph}

To connect the explicit tables to the short definition of $U$, write
$v_i$ for an accepted \emph{raw} vector and $r_i$ for its normalized
remainder. The indices in this subsection are the local indices of
\artifact{c8_reviewability_20260927/minimal_generation_witness.json}.
They must not be identified with the independently constructed
\texttt{r2\_000000} or \texttt{r3\_000000} checkpoint labels.

Here is a deterministic interpretation of every row of the next table.
Compute its raw vector $v_i$ using the indicated mode product on inputs
or earlier $r_j$. Select all earlier normalized rows of the same joint
Cartan weight. In increasing lexicographic order of their pivot states,
subtract the current coefficient at that pivot times its normalized
row. If the result is $b_i$, define
\[
 \pi_i=\min\operatorname{supp}(b_i),\qquad
 r_i=b_i/b_i[\pi_i].
\]
The state order first compares momentum tuples, then oscillator tuples.
The witness records all subtractions explicitly; the stated algorithm
specifies them without a large coefficient matrix. Distinct unit pivots
supply the independence proof. Since only earlier rows enter a recipe,
every $v_i$ and every $r_i$ belongs to the originally generated $U$.

The six current labels below are $H_L,E_L,F_L,H_R,E_R,F_R$.
The table gives the raw recipe, its Cartan weight, and its raw and
normalized support sizes. A support size counts lattice-oscillator
basis states, not rational components. The initial $24$ independent
rows are followed by $96$ accepted descendants. There are $644$
attempted insertions in the full queue, so rejected dependent or zero
candidates are not extra generators. The complete input conformal
vectors are also tabulated in Section~\ref{expdata:coordinates}.

\begingroup\small
\begin{longtable}{@{}rlcrr@{}}
\toprule $i$&Raw vector $v_i$&Weight&$|v_i|$&$|r_i|$\\\midrule\endfirsthead
\toprule $i$&Raw vector $v_i$&Weight&$|v_i|$&$|r_i|$\\\midrule\endhead
\bottomrule\endfoot
0 & $(H_L)_{-2}\mathbf1$ & $(0,0)$ & 6 & 6\\
1 & $(H_L)_{-1}H_L$ & $(0,0)$ & 21 & 21\\
2 & $(H_L)_{-1}E_L$ & $(2,0)$ & 6 & 6\\
3 & $(H_L)_{-1}F_L$ & $(-2,0)$ & 6 & 6\\
4 & $(H_L)_{-1}H_R$ & $(0,0)$ & 32 & 30\\
5 & $(H_L)_{-1}E_R$ & $(0,2)$ & 42 & 42\\
6 & $(H_L)_{-1}F_R$ & $(0,-2)$ & 42 & 42\\
7 & $(E_L)_{-1}H_R$ & $(2,0)$ & 7 & 6\\
8 & $(E_L)_{-1}E_R$ & $(2,2)$ & 7 & 7\\
9 & $(E_L)_{-1}F_R$ & $(2,-2)$ & 7 & 7\\
10 & $(F_L)_{-1}H_R$ & $(-2,0)$ & 7 & 6\\
11 & $(F_L)_{-1}E_R$ & $(-2,2)$ & 7 & 7\\
12 & $(F_L)_{-1}F_R$ & $(-2,-2)$ & 7 & 7\\
13 & $(H_R)_{-2}\mathbf1$ & $(0,0)$ & 7 & 6\\
14 & $(H_R)_{-1}H_R$ & $(0,0)$ & 28 & 21\\
15 & $(H_R)_{-1}E_R$ & $(0,2)$ & 49 & 42\\
16 & $(H_R)_{-1}F_R$ & $(0,-2)$ & 49 & 42\\
17 & $(E_R)_{-2}\mathbf1$ & $(0,2)$ & 7 & 7\\
18 & $(E_R)_{-1}E_R$ & $(0,4)$ & 26 & 26\\
19 & $(E_R)_{-1}F_R$ & $(0,0)$ & 46 & 60\\
20 & $(F_R)_{-2}\mathbf1$ & $(0,-2)$ & 7 & 42\\
21 & $(F_R)_{-1}F_R$ & $(0,-4)$ & 26 & 26\\
22 & $\omega_B$ & $(0,0)$ & 24 & 26\\
23 & $u$ & $(1,13)$ & 3 & 3\\
24 & $u_1r_{3}$ & $(-1,13)$ & 3 & 3\\
25 & $u_1r_{6}$ & $(1,11)$ & 7 & 7\\
26 & $u_1r_{12}$ & $(-1,11)$ & 7 & 7\\
27 & $u_1r_{21}$ & $(1,9)$ & 12 & 12\\
28 & $(F_R)_0r_{26}$ & $(-1,9)$ & 12 & 12\\
29 & $(F_R)_0r_{27}$ & $(1,7)$ & 22 & 22\\
30 & $(F_R)_0r_{28}$ & $(-1,7)$ & 22 & 22\\
31 & $(F_R)_0r_{29}$ & $(1,5)$ & 32 & 32\\
32 & $(F_R)_0r_{30}$ & $(-1,5)$ & 32 & 32\\
33 & $(F_R)_0r_{31}$ & $(1,3)$ & 39 & 39\\
34 & $(F_R)_0r_{32}$ & $(-1,3)$ & 39 & 39\\
35 & $(F_R)_0r_{33}$ & $(1,1)$ & 44 & 44\\
36 & $(F_R)_0r_{34}$ & $(-1,1)$ & 44 & 44\\
37 & $(F_R)_0r_{35}$ & $(1,-1)$ & 44 & 44\\
38 & $(F_R)_0r_{36}$ & $(-1,-1)$ & 44 & 44\\
39 & $(F_R)_0r_{37}$ & $(1,-3)$ & 39 & 39\\
40 & $(F_R)_0r_{38}$ & $(-1,-3)$ & 39 & 39\\
41 & $u_1r_{38}$ & $(0,12)$ & 9 & 9\\
42 & $(F_R)_0r_{39}$ & $(1,-5)$ & 32 & 32\\
43 & $(F_R)_0r_{40}$ & $(-1,-5)$ & 32 & 32\\
44 & $u_1r_{40}$ & $(0,10)$ & 16 & 16\\
45 & $(F_R)_0r_{42}$ & $(1,-7)$ & 22 & 22\\
46 & $(F_R)_0r_{43}$ & $(-1,-7)$ & 22 & 22\\
47 & $u_1r_{43}$ & $(0,8)$ & 32 & 32\\
48 & $(F_R)_0r_{44}$ & $(0,8)$ & 30 & 31\\
49 & $(F_R)_0r_{45}$ & $(1,-9)$ & 12 & 12\\
50 & $(F_R)_0r_{46}$ & $(-1,-9)$ & 12 & 12\\
51 & $u_1r_{46}$ & $(0,6)$ & 42 & 42\\
52 & $(F_R)_0r_{47}$ & $(0,6)$ & 45 & 44\\
53 & $(F_R)_0r_{49}$ & $(1,-11)$ & 7 & 7\\
54 & $(F_R)_0r_{50}$ & $(-1,-11)$ & 7 & 7\\
55 & $u_1r_{50}$ & $(0,4)$ & 36 & 42\\
56 & $(F_R)_0r_{51}$ & $(0,4)$ & 59 & 57\\
57 & $(F_R)_0r_{52}$ & $(0,4)$ & 58 & 57\\
58 & $(F_R)_0r_{53}$ & $(1,-13)$ & 3 & 3\\
59 & $(F_R)_0r_{54}$ & $(-1,-13)$ & 3 & 3\\
60 & $u_1r_{54}$ & $(0,2)$ & 35 & 30\\
61 & $(F_R)_0r_{55}$ & $(0,2)$ & 59 & 72\\
62 & $(F_R)_0r_{56}$ & $(0,2)$ & 67 & 70\\
63 & $u_1r_{59}$ & $(0,0)$ & 29 & 25\\
64 & $(F_R)_0r_{60}$ & $(0,0)$ & 64 & 62\\
65 & $(F_R)_0r_{61}$ & $(0,0)$ & 82 & 58\\
66 & $(F_R)_0r_{62}$ & $(0,0)$ & 82 & 63\\
67 & $(F_R)_0r_{63}$ & $(0,-2)$ & 49 & 49\\
68 & $(F_R)_0r_{64}$ & $(0,-2)$ & 72 & 68\\
69 & $(F_R)_0r_{65}$ & $(0,-2)$ & 70 & 71\\
70 & $(F_R)_0r_{67}$ & $(0,-4)$ & 42 & 42\\
71 & $(F_R)_0r_{68}$ & $(0,-4)$ & 61 & 58\\
72 & $(F_R)_0r_{69}$ & $(0,-4)$ & 61 & 57\\
73 & $(F_R)_0r_{70}$ & $(0,-6)$ & 42 & 42\\
74 & $(F_R)_0r_{71}$ & $(0,-6)$ & 45 & 44\\
75 & $(F_R)_0r_{73}$ & $(0,-8)$ & 32 & 32\\
76 & $(F_R)_0r_{74}$ & $(0,-8)$ & 32 & 31\\
77 & $u_1r_{74}$ & $(1,7)$ & 22 & 21\\
78 & $(F_R)_0r_{75}$ & $(0,-10)$ & 16 & 16\\
79 & $u_1r_{75}$ & $(1,5)$ & 32 & 31\\
80 & $u_1r_{76}$ & $(1,5)$ & 32 & 30\\
81 & $(F_L)_0r_{77}$ & $(-1,7)$ & 21 & 21\\
82 & $(F_R)_0r_{78}$ & $(0,-12)$ & 9 & 9\\
83 & $u_1r_{78}$ & $(1,3)$ & 30 & 38\\
84 & $(F_L)_0r_{79}$ & $(-1,5)$ & 31 & 31\\
85 & $(F_L)_0r_{80}$ & $(-1,5)$ & 30 & 30\\
86 & $(F_R)_0r_{80}$ & $(1,3)$ & 39 & 37\\
87 & $u_1r_{82}$ & $(1,1)$ & 27 & 42\\
88 & $(F_L)_0r_{83}$ & $(-1,3)$ & 38 & 38\\
89 & $(F_R)_0r_{83}$ & $(1,1)$ & 43 & 42\\
90 & $(F_R)_0r_{85}$ & $(-1,3)$ & 39 & 37\\
91 & $(F_R)_0r_{86}$ & $(1,1)$ & 43 & 41\\
92 & $(F_L)_0r_{87}$ & $(-1,1)$ & 42 & 42\\
93 & $(F_R)_0r_{87}$ & $(1,-1)$ & 44 & 43\\
94 & $(F_R)_0r_{88}$ & $(-1,1)$ & 43 & 42\\
95 & $(F_R)_0r_{89}$ & $(1,-1)$ & 44 & 42\\
96 & $(F_R)_0r_{90}$ & $(-1,1)$ & 43 & 41\\
97 & $(F_R)_0r_{91}$ & $(1,-1)$ & 43 & 41\\
98 & $(F_R)_0r_{92}$ & $(-1,-1)$ & 44 & 43\\
99 & $(F_R)_0r_{93}$ & $(1,-3)$ & 39 & 38\\
100 & $(F_R)_0r_{94}$ & $(-1,-1)$ & 44 & 42\\
101 & $(F_R)_0r_{95}$ & $(1,-3)$ & 39 & 37\\
102 & $(F_R)_0r_{96}$ & $(-1,-1)$ & 43 & 41\\
103 & $(F_R)_0r_{98}$ & $(-1,-3)$ & 39 & 38\\
104 & $(F_R)_0r_{99}$ & $(1,-5)$ & 31 & 31\\
105 & $(F_R)_0r_{100}$ & $(-1,-3)$ & 39 & 37\\
106 & $(F_R)_0r_{101}$ & $(1,-5)$ & 31 & 30\\
107 & $(F_R)_0r_{103}$ & $(-1,-5)$ & 31 & 31\\
108 & $(F_R)_0r_{104}$ & $(1,-7)$ & 22 & 21\\
109 & $(F_R)_0r_{105}$ & $(-1,-5)$ & 31 & 30\\
110 & $(F_R)_0r_{107}$ & $(-1,-7)$ & 22 & 21\\
111 & $u_1r_{107}$ & $(0,8)$ & 32 & 30\\
112 & $u_1r_{110}$ & $(0,6)$ & 40 & 43\\
113 & $(F_R)_0r_{112}$ & $(0,4)$ & 57 & 56\\
114 & $(F_R)_0r_{113}$ & $(0,2)$ & 66 & 71\\
115 & $(F_R)_0r_{114}$ & $(0,0)$ & 82 & 62\\
116 & $(F_R)_0r_{115}$ & $(0,-2)$ & 71 & 70\\
117 & $(F_R)_0r_{116}$ & $(0,-4)$ & 61 & 56\\
118 & $(F_R)_0r_{117}$ & $(0,-6)$ & 45 & 43\\
119 & $(F_R)_0r_{118}$ & $(0,-8)$ & 32 & 30\\
\end{longtable}
\endgroup
The recipe inventory has four accepted current derivatives, eighteen
accepted current normal products, $\omega_B$, $u$, seventy-eight
current-zero-mode descendants, and eighteen $u_1$ descendants. The
absence of separate accepted derivatives for $E_L$ and $F_L$ is a
linear-dependence outcome of the insertion order; it does not omit
those vectors from the initial affine span.

\subsubsection{Coordinates of all target fields in the generation basis}

For each target the next table lists all nonzero $c_i$ in
$\text{target}=\sum_i c_ir_i$. Together with the accepted recipe graph
and the normalization rule, these are explicit membership expressions
for every field in the preceding coordinate tables. They are often
much shorter than the fields' ambient supports. In particular the
three-term $u$ is exactly $r_{23}$, and the retained stress $e$ has
nine normalized coordinates.

\begingroup\small\renewcommand{\arraystretch}{1.4}
\begin{longtable}{@{}lrr@{}}
\toprule Target&$i$&Coefficient of $r_i$\\\midrule\endfirsthead
\toprule Target&$i$&Coefficient of $r_i$\\\midrule\endhead
\bottomrule\endfoot
$a_{0}$ & 18 & $-1$\\
 & 55 & $\frac{1159+589st}{10920}$\\
 & 56 & $\frac{38s-19t}{273}$\\
 & 57 & $\frac{-703s-171t}{1820}$\\
$a_{1}$ & 47 & $\frac{1-4st}{25}$\\
 & 48 & $\frac{-3579-2609st}{16920}$\\
 & 111 & $\frac{239283+48015s+218845t-1507st}{1602000}$\\
$a_{2}$ & 47 & $\frac{-7+3st}{20}$\\
 & 48 & $\frac{621+511st}{2256}$\\
 & 111 & $\frac{-11487-44835s+25795t+15323st}{213600}$\\
$a_{3}$ & 41 & $\frac{6+2s-4t+st}{14}$\\
$a_{4}$ & 35 & $\frac{217+32st}{295}$\\
 & 87 & $\frac{447s-319t}{1770}$\\
 & 89 & $\frac{-1953+1187st}{11800}$\\
 & 91 & $\frac{18717-2496s-12058t+377st}{143960}$\\
$a_{5}$ & 31 & $\frac{-55+15st}{126}$\\
 & 79 & $\frac{275-60s+120t-75st}{756}$\\
 & 80 & $\frac{-11565-30195s+13115t+6585st}{139608}$\\
$a_{6}$ & 29 & $\frac{307+297st}{1880}$\\
 & 77 & $\frac{-253s-319t}{1410}$\\
$a_{7}$ & 23 & $1$\\
$e$ & 1 & $\frac{890}{11439}$\\
 & 4 & $\frac{-860}{3813}$\\
 & 14 & $\frac{860}{3813}$\\
 & 19 & $\frac{260-205s}{3813}$\\
 & 22 & $\frac{40-15s}{93}$\\
 & 63 & $\frac{10s}{93}$\\
 & 64 & $\frac{10s}{93}$\\
 & 65 & $\frac{-11-st}{62}$\\
 & 66 & $\frac{23+3st}{62}$\\
$\omega_B$ & 1 & $\frac{2}{3}$\\
 & 4 & $-1$\\
 & 14 & $1$\\
 & 22 & $1$\\
\end{longtable}
\endgroup
The independent coordinate verifier reports $120$ normalization
identities, $120$ raw-generator expansions, and $20$ target expansions
(two bases for each of ten targets), together with the earlier-parent
and source-binding checks. This demonstrates the precise benefit of
saving a recipe graph: membership in the original generated algebra
can be reviewed separately from the linear equations in the ambient
basis. A coordinate verifier that merely sums saved raw vectors does
not by itself establish that those raw vectors were produced by the
named vertex modes; the exact exporter and mode-engine formulas supply
that remaining interpretation.

\subsection{Grade three: exact history and complete charge census}
\label{expfinite:grade-three}

The exact record \artifact{hi_cft_20260923_c8_grade3_exact_v2.json}
constructs $N_3=K_{-1}S_2+TS_2$ first, obtaining dimension $634$.
It then adds charged zero-products and closes under currents. Its
order differs from the original floating and finite-field runs, which
explains why the intermediate dimensions differ even though the final
space agrees. The full recorded exact progression is as follows.

\begin{longtable}{@{}lrr@{}}
\toprule Completed stage&Dimension&Cumulative seconds\\\midrule\endfirsthead
\toprule Completed stage&Dimension&Cumulative seconds\\\midrule\endhead
\bottomrule\endfoot
$N_3$ & 634 & 131.23\\
$(2, 14)$ charged family & 924 & 253.37\\
$(1, 13)$ charged family & 953 & 302.35\\
$(2, 8)$ charged family & 956 & 332.99\\
$(1, 9)$ charged family & 956 & 367.70\\
$(1, 9)$ charged family & 956 & 402.15\\
$(2, 6)$ charged family & 956 & 436.17\\
$(2, 2)$ charged family & 956 & 483.15\\
$(1, 5)$ charged family & 956 & 534.52\\
\end{longtable}
The quotient $P=S_3/N_3$ has dimension $322$. The next table
contains the complete saved charge census in a compressed Weyl-symmetric
form. Each row represents all independent sign choices of its nonzero
Cartan coordinates. Thus a row $(1,3)$ stands for four weight spaces,
a row $(0,4)$ for two, and $(0,0)$ for one. The orbit size is printed
explicitly. The $P$ column is transcribed by counting the quotient
weights in the exact saturation record, and the $N_3$ column is their
difference from the exact $S_3$ multiplicity.

\begin{longtable}{@{}crrrr@{}}
\toprule Weight representative&Sign orbit&$\dim S_3[w]$&$\dim N_3[w]$&$\dim P[w]$\\\midrule\endfirsthead
\toprule Weight representative&Sign orbit&$\dim S_3[w]$&$\dim N_3[w]$&$\dim P[w]$\\\midrule\endhead
\bottomrule\endfoot
$(0,0)$ & 1 & 58 & 40 & 18\\
$(0,2)$ & 2 & 49 & 31 & 18\\
$(0,4)$ & 2 & 37 & 22 & 15\\
$(0,6)$ & 2 & 31 & 17 & 14\\
$(0,8)$ & 2 & 22 & 13 & 9\\
$(0,10)$ & 2 & 14 & 7 & 7\\
$(0,12)$ & 2 & 8 & 4 & 4\\
$(0,14)$ & 2 & 4 & 1 & 3\\
$(0,16)$ & 2 & 1 & 0 & 1\\
$(0,18)$ & 2 & 1 & 0 & 1\\
$(1,1)$ & 4 & 29 & 19 & 10\\
$(1,3)$ & 4 & 25 & 16 & 9\\
$(1,5)$ & 4 & 22 & 14 & 8\\
$(1,7)$ & 4 & 16 & 10 & 6\\
$(1,9)$ & 4 & 10 & 6 & 4\\
$(1,11)$ & 4 & 7 & 5 & 2\\
$(1,13)$ & 4 & 5 & 4 & 1\\
$(1,15)$ & 4 & 1 & 1 & 0\\
$(2,0)$ & 2 & 12 & 12 & 0\\
$(2,2)$ & 4 & 7 & 7 & 0\\
$(2,4)$ & 4 & 5 & 5 & 0\\
$(2,6)$ & 4 & 3 & 3 & 0\\
$(2,8)$ & 4 & 3 & 3 & 0\\
$(2,10)$ & 4 & 1 & 1 & 0\\
$(2,12)$ & 4 & 1 & 1 & 0\\
\midrule Total (with sign orbits)& & $956$&$634$&$322$\\
\end{longtable}
The exact quotient current-character decomposition gives a second
way to check the $322$ count. Here $(j_L,j_R)$ denotes highest Cartan
weights, so the irreducible has dimension $(j_L+1)(j_R+1)$. The sum of
multiplicities is $28$. Equivalently, the quotient by the two raising
images has dimension $28$, and those images have rank $322-28=294$.

\begin{longtable}{@{}crrr@{}}
\toprule Highest weight&Multiplicity&Irrep dimension&Total dimension\\\midrule\endfirsthead
\toprule Highest weight&Multiplicity&Irrep dimension&Total dimension\\\midrule\endhead
\bottomrule\endfoot
$(1,13)$ & 1 & 28 & 28\\
$(1,11)$ & 1 & 24 & 24\\
$(1,9)$ & 2 & 20 & 40\\
$(1,7)$ & 2 & 16 & 32\\
$(1,5)$ & 2 & 12 & 24\\
$(1,3)$ & 1 & 8 & 8\\
$(1,1)$ & 1 & 4 & 4\\
$(0,18)$ & 1 & 19 & 19\\
$(0,14)$ & 2 & 15 & 30\\
$(0,12)$ & 1 & 13 & 13\\
$(0,10)$ & 3 & 11 & 33\\
$(0,8)$ & 2 & 9 & 18\\
$(0,6)$ & 5 & 7 & 35\\
$(0,4)$ & 1 & 5 & 5\\
$(0,2)$ & 3 & 3 & 9\\
\midrule Total&$28$&&$322$\\
\end{longtable}
The ambient weight-three space has dimension $11586$, so these
counts are dimensions of explicitly constructed subspaces of a much
larger space. Before the two closure gates, they do not rule out extra
vectors appearing from higher-weight products. After strong generation,
weight bookkeeping gives $U_3=S_3$, just as it gives $U_2=S_2$.

\subsubsection{What the preceding numerical and modular records establish}

The floating span record reached $634\to902\to956$ by using the
$(1,5)$ and then the first $(1,9)$ families. Its smallest accepted
relative residual was $9.7730725\times10^{-5}$; its largest rejected
relative residual was $2.6155948\times10^{-9}$. These provide evidence
of numerical separation for that run, not an exact upper bound.

The modular run used $p=1000003$ with
$s\mapsto231543$ and $t\mapsto1000$. At a prime where all required
denominators remain invertible, a nonzero minor survives as a proof of
a characteristic-zero lower bound. A modular rank of $956$ therefore
shows that at least $956$ vectors are independent. It does not show
that a characteristic-zero span has no further vectors. The separate
number-field run is what records exact dimension $956$.

The historical single-generator action test used only the three-term
$u$, on all $956$ numerical source vectors. Its diagnostics are:

\begin{longtable}{@{}crrr@{}}
\toprule Mode&Target grade&Total leakage&Largest block relative leakage\\\midrule\endfirsthead
\toprule Mode&Target grade&Total leakage&Largest block relative leakage\\\midrule\endhead
\bottomrule\endfoot
1 & 3 & $4.111\times10^{-10}$ & $6.343\times10^{-11}$\\
2 & 2 & $1.524\times10^{-10}$ & $5.408\times10^{-11}$\\
3 & 1 & $6.015\times10^{-11}$ & $5.237\times10^{-11}$\\
\end{longtable}
These small floating residuals motivated the exact closure work.
They do not establish the action of every composite generator, and they
do not replace the finite gates. The exact positive-mode consequences
proved after Gate~4 apply to the whole $S_2$ and therefore have a
strictly different logical scope from this early single-generator test.


\subsection{The 3970 lower identities: an explicit task inventory}
\label{expfinite:3970}

The number $3970$ in the closure ledger is the size of a concrete
ordered task list. It is useful to unpack that number because an
unqualified ``all lower checks passed'' conceals several different
target spaces. The implementation
\artifact{hi_cft_20260923_next_structure_lower_closure.py} uses ten
weight-two primitive fields: the eight $a_i$, followed by $e,f$.
Its tasks are exactly:
\begin{center}\small
\begin{tabular}{@{}lrr@{}}
\toprule
Task family&Number&Target grade\\\midrule
$L(1)a=0$ for the ten primitive fields&$10$&$1$\\
$a_0v\in S_3$ for ten fields and $120$ basis vectors&$1200$&$3$\\
$a_1v\in S_2$ for the same pairs&$1200$&$2$\\
$x_0v\in S_2$ for six currents and $120$ basis vectors&$720$&$2$\\
$x_1v\in K$ for the same pairs&$720$&$1$\\
$L(1)v\in K$ for all $120$ basis vectors&$120$&$1$\\
\midrule Total&$3970$&\\\bottomrule
\end{tabular}
\end{center}
The first two left fields processed in the weight-two product portion
are $e,f$, since their zero-product families were not seeds in the old
grade-three construction. The source then processes $a_0,\ldots,a_7$.
This order is why database primary keys cannot be interpreted by simply
multiplying the charged-field index by $120$.

There is a further basis convention in this database. The six
\texttt{current\_basis} rows are unit-pivot echelon rows, and are
\[
\begin{array}{c|l}
0&-H_L\\
1&E_L\\
2&-F_L\\
3&-(H_R+5H_L)/2\\
4&E_R\\
5&-F_R/5
\end{array}
\]
in the canonical current notation. The left current in a stored
\texttt{current} task is still the original canonical current indexed
$0,\ldots,5$; the transformed rows appear when recording the target's
coordinates in $K$. Interchanging these two roles changes coefficients
and can appear to create a sign failure even when the exact identity
is correct.

The completed record has cursor $3970$, all residual lists empty,
$224480$ sparse input terms, and a largest recorded coefficient size
of $290$ bits. Its exact arithmetic uses four GMP rational components.
Its saved elapsed time is $77.88$ seconds and peak RSS $239784$ KiB.
These resource figures describe the old run; they are not estimates for
arbitrarily larger grades.

\subsection{The actual triplet before any quotient-dimension assertion}
\label{expfinite:actual-triplet}

The triplet record
\artifact{hi_cft_20260923_next_structure_p4_controls.json}
starts with the actual vector
\[
 t=(a_0)_0\texttt{r3\_000705}\in B_4,
 \qquad \operatorname{wt}_{\rm Cartan}(t)=(0,2).
\]
It constructs its highest triplet vector by
\begin{align*}
 v={}&\left(1-\frac{F_LE_L}{2}\right)
       \left(1-\frac{F_LE_L}{6}\right)
       \prod_{d\in D_2}\left(1-\frac{F_RE_R}{d}\right)t,\\
 D_2={}&\{4,10,18,28,40,54,70,88,108,130,154\}.
\end{align*}
Every operator here is a zero mode, with $E$ applied before $F$.
The three stored fields are
\[
 P4_{+}=v,\qquad P4_0=(F_R)_0v,\qquad
 P4_{-}=(F_R)_0^2v.
\]
They have $2313$, $2363$, and $2313$ nonzero lattice-oscillator
coordinates, respectively. The left projection retains $2669$ terms;
the first right projection lowers that support to $2313$, and the
remaining right factors retain this count. Support counts are a useful
fingerprint for a reconstruction, but are not an equality test.

There are eighteen current-zero-mode identities, six for each of the
three actual vectors. They establish the current type $(1,3)$ in
representation-dimension notation. Six further exact checks put
$L(1)P4_{\pm,0}$ in $S_3$ and $L(2)P4_{\pm,0}$ in $S_2$.
These direct positive-Virasoro controls precede Gate~4 in the logical
order, so subsequent positive-mode arguments do not use their own
conclusion as a premise.

The separate modular replacement control finds nonzero multipliers
$779530$, $591896$, and $665783$ at $p=1000003$ for the three
weights $(0,2),(0,0),(0,-2)$. It shows that these actual fields replace
the selected discovery columns. It does not prove that $N_4$ has a
particular characteristic-zero dimension. The later argument therefore
allows both possibilities: $P_4\subset N_4$, or an irreducible triplet
survives in $W_4/N_4$.

\subsection{Gate 4 saturation: every operator deficit and every added input}
\label{expfinite:gate4-inventory}

The exact saturation record contains $644$ raising-coordinate columns,
two for each basis vector of $P=S_3/N_3$. Their rank is $294$.
For each highest left input, the old exact seed inputs are appended to
these columns. The next table displays the resulting rank and the exact
deficit. The index $8$ denotes the retained neutral stress $e$; indices
$0,\ldots,7$ denote the charged highest fields. Each augmented rank
becomes $322$ after the listed additional requests are supplied.

\begin{longtable}{@{}rrrrr@{}}
\toprule Left index&Old seed inputs&Rank with raising images&Deficit&Completed rank\\\midrule\endfirsthead
\toprule Left index&Old seed inputs&Rank with raising images&Deficit&Completed rank\\\midrule\endhead
\bottomrule\endfoot
0 & 197 & 307 & 15 & 322\\
1 & 236 & 303 & 19 & 322\\
2 & 234 & 302 & 20 & 322\\
3 & 277 & 306 & 16 & 322\\
4 & 226 & 313 & 9 & 322\\
5 & 244 & 307 & 15 & 322\\
6 & 259 & 305 & 17 & 322\\
7 & 302 & 313 & 9 & 322\\
8 & 170 & 312 & 10 & 322\\
\midrule Total seed inputs&$2145$&&$130$&\\
\end{longtable}
The $130$ requested input pairs are printed next. ``Direct'' means
one of the $118$ completed additional ordinary exact identities.
``Singlet'' means one of the twelve weight-zero targets recovered by
raising controls and the projected exact identity. These are descriptions
of the completed later coverage, not success claims made by the original
request file. The source list is
\artifact{hi_cft_20260923_c8two_structure_saturation.json}; final
coverage is recorded in
\artifact{hi_cft_20260923_c8two_gate4_verified.json}.

\begingroup\small
\begin{longtable}{@{}rlccl@{}}
\toprule Left&Exact right input&Input weight&Output weight&Coverage\\\midrule\endfirsthead
\toprule Left&Exact right input&Input weight&Output weight&Coverage\\\midrule\endhead
\bottomrule\endfoot
0 & \texttt{r3\_000921} & $(0,-10)$ & $(0,-6)$ & Direct\\
0 & \texttt{r3\_000917} & $(0,-10)$ & $(0,-6)$ & Direct\\
0 & \texttt{r3\_000915} & $(0,-10)$ & $(0,-6)$ & Direct\\
0 & \texttt{r3\_000920} & $(0,-8)$ & $(0,-4)$ & Direct\\
0 & \texttt{r3\_000916} & $(0,-8)$ & $(0,-4)$ & Direct\\
0 & \texttt{r3\_000874} & $(-1,-3)$ & $(-1,1)$ & Direct\\
0 & \texttt{r3\_000952} & $(0,-6)$ & $(0,-2)$ & Direct\\
0 & \texttt{r3\_000951} & $(0,-6)$ & $(0,-2)$ & Direct\\
0 & \texttt{r3\_000946} & $(0,-6)$ & $(0,-2)$ & Direct\\
0 & \texttt{r3\_000913} & $(0,-6)$ & $(0,-2)$ & Direct\\
0 & \texttt{r3\_000912} & $(0,-6)$ & $(0,-2)$ & Direct\\
0 & \texttt{r3\_000955} & $(0,-2)$ & $(0,2)$ & Direct\\
0 & \texttt{r3\_000947} & $(0,-2)$ & $(0,2)$ & Direct\\
0 & \texttt{r3\_000945} & $(0,-2)$ & $(0,2)$ & Direct\\
0 & \texttt{r3\_000950} & $(0,-4)$ & $(0,0)$ & Singlet\\
1 & \texttt{r3\_000923} & $(0,-14)$ & $(0,-6)$ & Direct\\
1 & \texttt{r3\_000910} & $(0,-14)$ & $(0,-6)$ & Direct\\
1 & \texttt{r3\_000955} & $(0,-2)$ & $(0,6)$ & Direct\\
1 & \texttt{r3\_000947} & $(0,-2)$ & $(0,6)$ & Direct\\
1 & \texttt{r3\_000945} & $(0,-2)$ & $(0,6)$ & Direct\\
1 & \texttt{r3\_000922} & $(0,-12)$ & $(0,-4)$ & Direct\\
1 & \texttt{r3\_000950} & $(0,-4)$ & $(0,4)$ & Direct\\
1 & \texttt{r3\_000921} & $(0,-10)$ & $(0,-2)$ & Direct\\
1 & \texttt{r3\_000917} & $(0,-10)$ & $(0,-2)$ & Direct\\
1 & \texttt{r3\_000915} & $(0,-10)$ & $(0,-2)$ & Direct\\
1 & \texttt{r3\_000898} & $(-1,-7)$ & $(-1,1)$ & Direct\\
1 & \texttt{r3\_000882} & $(-1,-7)$ & $(-1,1)$ & Direct\\
1 & \texttt{r3\_000952} & $(0,-6)$ & $(0,2)$ & Direct\\
1 & \texttt{r3\_000951} & $(0,-6)$ & $(0,2)$ & Direct\\
1 & \texttt{r3\_000946} & $(0,-6)$ & $(0,2)$ & Direct\\
1 & \texttt{r3\_000913} & $(0,-6)$ & $(0,2)$ & Direct\\
1 & \texttt{r3\_000912} & $(0,-6)$ & $(0,2)$ & Direct\\
1 & \texttt{r3\_000920} & $(0,-8)$ & $(0,0)$ & Singlet\\
1 & \texttt{r3\_000916} & $(0,-8)$ & $(0,0)$ & Singlet\\
2 & \texttt{r3\_000923} & $(0,-14)$ & $(0,-6)$ & Direct\\
2 & \texttt{r3\_000910} & $(0,-14)$ & $(0,-6)$ & Direct\\
2 & \texttt{r3\_000955} & $(0,-2)$ & $(0,6)$ & Direct\\
2 & \texttt{r3\_000947} & $(0,-2)$ & $(0,6)$ & Direct\\
2 & \texttt{r3\_000945} & $(0,-2)$ & $(0,6)$ & Direct\\
2 & \texttt{r3\_000922} & $(0,-12)$ & $(0,-4)$ & Direct\\
2 & \texttt{r3\_000950} & $(0,-4)$ & $(0,4)$ & Direct\\
2 & \texttt{r3\_000900} & $(-1,-9)$ & $(-1,-1)$ & Direct\\
2 & \texttt{r3\_000921} & $(0,-10)$ & $(0,-2)$ & Direct\\
2 & \texttt{r3\_000917} & $(0,-10)$ & $(0,-2)$ & Direct\\
2 & \texttt{r3\_000915} & $(0,-10)$ & $(0,-2)$ & Direct\\
2 & \texttt{r3\_000898} & $(-1,-7)$ & $(-1,1)$ & Direct\\
2 & \texttt{r3\_000882} & $(-1,-7)$ & $(-1,1)$ & Direct\\
2 & \texttt{r3\_000952} & $(0,-6)$ & $(0,2)$ & Direct\\
2 & \texttt{r3\_000951} & $(0,-6)$ & $(0,2)$ & Direct\\
2 & \texttt{r3\_000946} & $(0,-6)$ & $(0,2)$ & Direct\\
2 & \texttt{r3\_000913} & $(0,-6)$ & $(0,2)$ & Direct\\
2 & \texttt{r3\_000912} & $(0,-6)$ & $(0,2)$ & Direct\\
2 & \texttt{r3\_000920} & $(0,-8)$ & $(0,0)$ & Singlet\\
2 & \texttt{r3\_000916} & $(0,-8)$ & $(0,0)$ & Singlet\\
3 & \texttt{r3\_000919} & $(0,-18)$ & $(0,-6)$ & Direct\\
3 & \texttt{r3\_000952} & $(0,-6)$ & $(0,6)$ & Direct\\
3 & \texttt{r3\_000951} & $(0,-6)$ & $(0,6)$ & Direct\\
3 & \texttt{r3\_000946} & $(0,-6)$ & $(0,6)$ & Direct\\
3 & \texttt{r3\_000913} & $(0,-6)$ & $(0,6)$ & Direct\\
3 & \texttt{r3\_000912} & $(0,-6)$ & $(0,6)$ & Direct\\
3 & \texttt{r3\_000920} & $(0,-8)$ & $(0,4)$ & Direct\\
3 & \texttt{r3\_000916} & $(0,-8)$ & $(0,4)$ & Direct\\
3 & \texttt{r3\_000902} & $(-1,-13)$ & $(-1,-1)$ & Direct\\
3 & \texttt{r3\_000923} & $(0,-14)$ & $(0,-2)$ & Direct\\
3 & \texttt{r3\_000910} & $(0,-14)$ & $(0,-2)$ & Direct\\
3 & \texttt{r3\_000901} & $(-1,-11)$ & $(-1,1)$ & Direct\\
3 & \texttt{r3\_000921} & $(0,-10)$ & $(0,2)$ & Direct\\
3 & \texttt{r3\_000917} & $(0,-10)$ & $(0,2)$ & Direct\\
3 & \texttt{r3\_000915} & $(0,-10)$ & $(0,2)$ & Direct\\
3 & \texttt{r3\_000922} & $(0,-12)$ & $(0,0)$ & Singlet\\
4 & \texttt{r3\_000898} & $(-1,-7)$ & $(0,-6)$ & Direct\\
4 & \texttt{r3\_000882} & $(-1,-7)$ & $(0,-6)$ & Direct\\
4 & \texttt{r3\_000897} & $(-1,-5)$ & $(0,-4)$ & Direct\\
4 & \texttt{r3\_000879} & $(-1,-5)$ & $(0,-4)$ & Direct\\
4 & \texttt{r3\_000874} & $(-1,-3)$ & $(0,-2)$ & Direct\\
4 & \texttt{r3\_000955} & $(0,-2)$ & $(1,-1)$ & Direct\\
4 & \texttt{r3\_000947} & $(0,-2)$ & $(1,-1)$ & Direct\\
4 & \texttt{r3\_000945} & $(0,-2)$ & $(1,-1)$ & Direct\\
4 & \texttt{r3\_000873} & $(-1,-1)$ & $(0,0)$ & Singlet\\
5 & \texttt{r3\_000901} & $(-1,-11)$ & $(0,-6)$ & Direct\\
5 & \texttt{r3\_000900} & $(-1,-9)$ & $(0,-4)$ & Direct\\
5 & \texttt{r3\_000899} & $(-1,-9)$ & $(0,-4)$ & Direct\\
5 & \texttt{r3\_000873} & $(-1,-1)$ & $(0,4)$ & Direct\\
5 & \texttt{r3\_000898} & $(-1,-7)$ & $(0,-2)$ & Direct\\
5 & \texttt{r3\_000882} & $(-1,-7)$ & $(0,-2)$ & Direct\\
5 & \texttt{r3\_000874} & $(-1,-3)$ & $(0,2)$ & Direct\\
5 & \texttt{r3\_000952} & $(0,-6)$ & $(1,-1)$ & Direct\\
5 & \texttt{r3\_000951} & $(0,-6)$ & $(1,-1)$ & Direct\\
5 & \texttt{r3\_000946} & $(0,-6)$ & $(1,-1)$ & Direct\\
5 & \texttt{r3\_000913} & $(0,-6)$ & $(1,-1)$ & Direct\\
5 & \texttt{r3\_000912} & $(0,-6)$ & $(1,-1)$ & Direct\\
5 & \texttt{r3\_000950} & $(0,-4)$ & $(1,1)$ & Direct\\
5 & \texttt{r3\_000897} & $(-1,-5)$ & $(0,0)$ & Singlet\\
5 & \texttt{r3\_000879} & $(-1,-5)$ & $(0,0)$ & Singlet\\
6 & \texttt{r3\_000902} & $(-1,-13)$ & $(0,-6)$ & Direct\\
6 & \texttt{r3\_000873} & $(-1,-1)$ & $(0,6)$ & Direct\\
6 & \texttt{r3\_000901} & $(-1,-11)$ & $(0,-4)$ & Direct\\
6 & \texttt{r3\_000874} & $(-1,-3)$ & $(0,4)$ & Direct\\
6 & \texttt{r3\_000900} & $(-1,-9)$ & $(0,-2)$ & Direct\\
6 & \texttt{r3\_000899} & $(-1,-9)$ & $(0,-2)$ & Direct\\
6 & \texttt{r3\_000897} & $(-1,-5)$ & $(0,2)$ & Direct\\
6 & \texttt{r3\_000920} & $(0,-8)$ & $(1,-1)$ & Direct\\
6 & \texttt{r3\_000879} & $(-1,-5)$ & $(0,2)$ & Direct\\
6 & \texttt{r3\_000916} & $(0,-8)$ & $(1,-1)$ & Direct\\
6 & \texttt{r3\_000952} & $(0,-6)$ & $(1,1)$ & Direct\\
6 & \texttt{r3\_000951} & $(0,-6)$ & $(1,1)$ & Direct\\
6 & \texttt{r3\_000946} & $(0,-6)$ & $(1,1)$ & Direct\\
6 & \texttt{r3\_000913} & $(0,-6)$ & $(1,1)$ & Direct\\
6 & \texttt{r3\_000912} & $(0,-6)$ & $(1,1)$ & Direct\\
6 & \texttt{r3\_000898} & $(-1,-7)$ & $(0,0)$ & Singlet\\
6 & \texttt{r3\_000882} & $(-1,-7)$ & $(0,0)$ & Singlet\\
7 & \texttt{r3\_000898} & $(-1,-7)$ & $(0,6)$ & Direct\\
7 & \texttt{r3\_000882} & $(-1,-7)$ & $(0,6)$ & Direct\\
7 & \texttt{r3\_000900} & $(-1,-9)$ & $(0,4)$ & Direct\\
7 & \texttt{r3\_000899} & $(-1,-9)$ & $(0,4)$ & Direct\\
7 & \texttt{r3\_000923} & $(0,-14)$ & $(1,-1)$ & Direct\\
7 & \texttt{r3\_000910} & $(0,-14)$ & $(1,-1)$ & Direct\\
7 & \texttt{r3\_000901} & $(-1,-11)$ & $(0,2)$ & Direct\\
7 & \texttt{r3\_000922} & $(0,-12)$ & $(1,1)$ & Direct\\
7 & \texttt{r3\_000902} & $(-1,-13)$ & $(0,0)$ & Singlet\\
8 & \texttt{r3\_000952} & $(0,-6)$ & $(0,-6)$ & Direct\\
8 & \texttt{r3\_000951} & $(0,-6)$ & $(0,-6)$ & Direct\\
8 & \texttt{r3\_000946} & $(0,-6)$ & $(0,-6)$ & Direct\\
8 & \texttt{r3\_000913} & $(0,-6)$ & $(0,-6)$ & Direct\\
8 & \texttt{r3\_000912} & $(0,-6)$ & $(0,-6)$ & Direct\\
8 & \texttt{r3\_000950} & $(0,-4)$ & $(0,-4)$ & Direct\\
8 & \texttt{r3\_000873} & $(-1,-1)$ & $(-1,-1)$ & Direct\\
8 & \texttt{r3\_000955} & $(0,-2)$ & $(0,-2)$ & Direct\\
8 & \texttt{r3\_000947} & $(0,-2)$ & $(0,-2)$ & Direct\\
8 & \texttt{r3\_000945} & $(0,-2)$ & $(0,-2)$ & Direct\\
\end{longtable}
\endgroup
There is no assertion that these seed inputs are stable under both
lowering operators, or that they exhaust all $322$ quotient basis
vectors for a fixed left field. The deduction uses stability of the
\emph{kernel} of $[p]\mapsto[a_0p]$ under the raising operators.
For a highest $a$, covariance makes that stability immediate. If the
raising images and certified kernel inputs span $P$, then the quotient
by this kernel equals its own image under the raising ideal. Because
every raising word of length $25$ vanishes on grade three, iteration
forces the quotient to vanish. This explains why the table's nine
completed rank-$322$ assertions imply $9\cdot322=2898$ highest-left
identities, despite far fewer direct computations for each family.

\subsubsection{The twelve additional raising-premise inputs}

The singlet route also uses twelve ordinary identities with output
Cartan weight $(0,2)$. These belong to a separate request file and should
not be silently identified with the twelve projected weight-zero targets.
The following table gives the complete list from
\artifact{hi_cft_20260923_c8two_structure_h4_raise_request.json}.
The request file itself has status ``request only''; the completed
raised-input summary and singlet-premise verifier are the later proof
records.

\begin{longtable}{@{}rlcc@{}}
\toprule Left&Exact right input&Input weight&Output weight\\\midrule\endfirsthead
\toprule Left&Exact right input&Input weight&Output weight\\\midrule\endhead
\bottomrule\endfoot
0 & \texttt{r3\_000944} & $(0,-2)$ & $(0,2)$\\
1 & \texttt{r3\_000905} & $(0,-6)$ & $(0,2)$\\
1 & \texttt{r3\_000907} & $(0,-6)$ & $(0,2)$\\
2 & \texttt{r3\_000905} & $(0,-6)$ & $(0,2)$\\
2 & \texttt{r3\_000907} & $(0,-6)$ & $(0,2)$\\
3 & \texttt{r3\_000909} & $(0,-10)$ & $(0,2)$\\
4 & \texttt{r3\_000842} & $(-1,1)$ & $(0,2)$\\
5 & \texttt{r3\_000850} & $(-1,-3)$ & $(0,2)$\\
5 & \texttt{r3\_000871} & $(-1,-3)$ & $(0,2)$\\
6 & \texttt{r3\_000854} & $(-1,-5)$ & $(0,2)$\\
6 & \texttt{r3\_000875} & $(-1,-5)$ & $(0,2)$\\
7 & \texttt{r3\_000885} & $(-1,-11)$ & $(0,2)$\\
\end{longtable}
Together with the earlier exact seed relations, these imply the
$24$ raising conditions $E_Lt,E_Rt\in W_4$ for the twelve remaining
weight-zero targets. The source-bound replay is
\artifact{hi_cft_20260923_c8two_h4_singlet_premise_verified.json}.
Only after this premise is available does a projected residual prove an
unprojected membership statement.

\subsection{All twenty-one normal columns in the singlet certificate}
\label{expfinite:projected-columns}

For weight-zero inputs the joint-singlet projector has the right-axis
denominators
\[
D_0=\{2,6,12,20,30,42,56,72,90,110,132,156\},
\]
with the same left-axis denominators $2,6$. Thus $D_0$ differs from
$D_2$ above: they project different Cartan weight spaces. The exact
projected equations have the form
\[
 \Pi t_i=\sum_{j=0}^{20}c_{ij}\Pi n_j.
\]
The following table exposes every unprojected word $n_j$ and its
original normal-column ID. The projected-column number is a separate
index. This translation is enough to interpret the saved coefficients
as actual elements of $N_4$, rather than as arbitrary vectors in a
linear solver.

\begin{longtable}{@{}rrl@{}}
\toprule Projected $j$&Original normal ID&Actual word $n_j$\\\midrule\endfirsthead
\toprule Projected $j$&Original normal ID&Actual word $n_j$\\\midrule\endhead
\bottomrule\endfoot
0 & 1 & $(H_L)_{-1}$\texttt{r3\_000001}\\
1 & 14 & $(H_L)_{-1}$\texttt{r3\_000014}\\
2 & 22 & $(H_L)_{-1}$\texttt{r3\_000022}\\
3 & 23 & $(H_L)_{-1}$\texttt{r3\_000023}\\
4 & 121 & $(H_L)_{-1}$\texttt{r3\_000121}\\
5 & 1811 & $(H_R)_{-1}$\texttt{r3\_000221}\\
6 & 1816 & $(H_R)_{-1}$\texttt{r3\_000226}\\
7 & 1819 & $(H_R)_{-1}$\texttt{r3\_000229}\\
8 & 1820 & $(H_R)_{-1}$\texttt{r3\_000230}\\
9 & 1823 & $(H_R)_{-1}$\texttt{r3\_000233}\\
10 & 1920 & $(H_R)_{-1}$\texttt{r3\_000330}\\
11 & 2256 & $(H_R)_{-1}$\texttt{r3\_000666}\\
12 & 2521 & $(H_R)_{-1}$\texttt{r3\_000931}\\
13 & 2544 & $(H_R)_{-1}$\texttt{r3\_000954}\\
14 & 4003 & $T$\texttt{r3\_000536}\\
15 & 4004 & $T$\texttt{r3\_000537}\\
16 & 4429 & $(a_{0})_{-1}$\texttt{r2\_000028}\\
17 & 4535 & $(a_{1})_{-1}$\texttt{r2\_000037}\\
18 & 4544 & $(a_{1})_{-1}$\texttt{r2\_000046}\\
19 & 4640 & $(a_{2})_{-1}$\texttt{r2\_000046}\\
20 & 4735 & $(a_{3})_{-1}$\texttt{r2\_000059}\\
\end{longtable}
The exact projected coefficients can be long: some have numerators
and denominators of approximately one hundred decimal digits. A compact
exact description of their reconstruction is to choose $21$ independent
pivot coordinates of the projected normal columns, form the square
matrix $M$ of their coefficients, and solve
\[
 c_i=M^{-1}(\Pi t_i)|_{\rm pivots}.
\]
This produces candidate coefficients only. The decisive check is the
entire sparse identity $\Pi t_i-\sum_jc_{ij}\Pi n_j=0$, on every
ambient basis state and all four rational components. A pivot solve
without that last check could pass while a nonzero residual remains
outside the chosen coordinates.

The following target-size table is transcribed from the portable
independent replay. All twelve full residuals were zero.

\begin{longtable}{@{}rrlrr@{}}
\toprule Target ID&Left&Right input&Target terms&Columns used\\\midrule\endfirsthead
\toprule Target ID&Left&Right input&Target terms&Columns used\\\midrule\endhead
\bottomrule\endfoot
0 & 0 & \texttt{r3\_000950} & 2343 & 19\\
1 & 1 & \texttt{r3\_000916} & 2343 & 19\\
2 & 1 & \texttt{r3\_000920} & 2343 & 19\\
3 & 2 & \texttt{r3\_000916} & 2343 & 19\\
4 & 2 & \texttt{r3\_000920} & 2343 & 19\\
5 & 3 & \texttt{r3\_000922} & 2343 & 19\\
6 & 4 & \texttt{r3\_000873} & 2726 & 21\\
7 & 5 & \texttt{r3\_000879} & 2726 & 21\\
8 & 5 & \texttt{r3\_000897} & 2726 & 21\\
9 & 6 & \texttt{r3\_000882} & 2726 & 21\\
10 & 6 & \texttt{r3\_000898} & 2726 & 21\\
11 & 7 & \texttt{r3\_000902} & 2726 & 21\\
\end{longtable}
The portable packet contains the twelve projected targets and the
$21$ projected normal columns, with $66610$ sparse state terms in
all. It permits exact residual checking without the original SQLite
stores. Its historical \texttt{grade5:} prefix for projected normal
vectors is a storage name; the vectors in this packet have weight four.
The portable residual replay does not reconstruct the projector action,
the normal-word provenance, or the raising premises, all of which remain
separately source-bound inputs.

The final recovery is a quotient calculation. Since $E_Lt_i,E_Rt_i$
lie in $W_4$ and $W_4$ is current-stable, every $FE$ factor acts as
zero on $[t_i]$. Hence $[\Pi t_i]=[t_i]$ in $B_4/W_4$.
Every $\Pi n_j$ stays in the current-stable space $N_4$. Therefore
the full residual equality implies $[t_i]=0$. The argument never needs
the $21$ selected columns to span all of $N_4$.


\subsection{Gate 5: the charged-span witness behind the neutral step}
\label{expfinite:gate5-seeds}

The charged part of Gate~5 is unusually small after Gate~4 is known.
The eight charged highest weights are
\[
(0,4),(0,8),(0,8),(0,12),(1,1),(1,5),(1,7),(1,13).
\]
If $Q=W_4/N_4$ is nonzero, it is the triplet with weights
$(0,-2),(0,0),(0,2)$. Two previously computed bottom relations give
nonzero coefficients of its lowest vector. Seven of the eight required
charged brackets vanish exactly, and the remaining bracket is an
explicit current normal product. The complete small coefficients and
normal-word identities are printed in the accompanying finite-witness
section. They imply $a_0P_4\subset N_5$ for every charged $a$.

A separate spanning witness is required for the retained neutral field.
This subsection exposes the $28$ accepted seeds of that witness.
For a checkpoint row $r$ of $S_2$, the seed indexed by $(i,x,r)$ is
\[
       \bigl[(a_i)_0(x_0r)\bigr]\in P=S_3/N_3.
\]
The saved lower identities express $x_0r$ in the exact $S_2$ basis,
and express every $(a_i)_0r'$ in the exact $S_3$ basis. Composing
these two saved coordinate maps gives the displayed seed. Thus no new
grade-three mode calculation is needed for this particular spanning
replay. The seed lies in the current span of charged brackets: a
current zero mode kills the neutral singlets, while the affine-normal
right parts contribute zero modulo $N_3$ by the lower normal controls.

Before inserting these seeds, the $644$ lowering images of the $322$
quotient basis vectors have rank $294$. Each row below adds one to
that rank; the final rank is $322$. The original producer tested
$1129$ candidates before obtaining this accepted list.

\begin{longtable}{@{}rrllr@{}}
\toprule Seed&$i$&Current $x$&Right input $r$&Rank after insertion\\\midrule\endfirsthead
\toprule Seed&$i$&Current $x$&Right input $r$&Rank after insertion\\\midrule\endhead
\bottomrule\endfoot
0 & 0 & $H_L$ & \texttt{r2\_000060} & 295\\
1 & 0 & $H_L$ & \texttt{r2\_000061} & 296\\
2 & 0 & $H_L$ & \texttt{r2\_000064} & 297\\
3 & 0 & $H_L$ & \texttt{r2\_000065} & 298\\
4 & 0 & $H_L$ & \texttt{r2\_000066} & 299\\
5 & 0 & $H_L$ & \texttt{r2\_000068} & 300\\
6 & 0 & $H_L$ & \texttt{r2\_000076} & 301\\
7 & 0 & $H_L$ & \texttt{r2\_000077} & 302\\
8 & 0 & $H_L$ & \texttt{r2\_000078} & 303\\
9 & 0 & $H_L$ & \texttt{r2\_000094} & 304\\
10 & 0 & $H_R$ & \texttt{r2\_000025} & 305\\
11 & 0 & $H_R$ & \texttt{r2\_000027} & 306\\
12 & 0 & $H_R$ & \texttt{r2\_000029} & 307\\
13 & 0 & $H_R$ & \texttt{r2\_000030} & 308\\
14 & 0 & $H_R$ & \texttt{r2\_000031} & 309\\
15 & 0 & $H_R$ & \texttt{r2\_000032} & 310\\
16 & 0 & $H_R$ & \texttt{r2\_000039} & 311\\
17 & 0 & $H_R$ & \texttt{r2\_000040} & 312\\
18 & 0 & $H_R$ & \texttt{r2\_000041} & 313\\
19 & 0 & $H_R$ & \texttt{r2\_000048} & 314\\
20 & 0 & $H_R$ & \texttt{r2\_000050} & 315\\
21 & 0 & $E_R$ & \texttt{r2\_000034} & 316\\
22 & 1 & $H_R$ & \texttt{r2\_000030} & 317\\
23 & 1 & $H_R$ & \texttt{r2\_000034} & 318\\
24 & 1 & $H_R$ & \texttt{r2\_000036} & 319\\
25 & 1 & $H_R$ & \texttt{r2\_000043} & 320\\
26 & 1 & $H_R$ & \texttt{r2\_000045} & 321\\
27 & 1 & $H_R$ & \texttt{r2\_000048} & 322\\
\end{longtable}
These columns are specialized at the same good split prime
$p=1000003$, with $s\mapsto231543$, $t\mapsto1000$.
The independent verifier reconstructs the seeds from the exact lower
relations, checks the two $\mathfrak{sl}_2$ commutators and the
cross-axis commutators against the saved raising maps, and recomputes
the positive full-rank witness. In particular, the lowering maps are
not merely arbitrary matrices with the desired rank.

Let $C$ be the full current submodule generated by charged brackets.
The rank witness gives
\[
 P=F_LP+F_RP+C.
\]
On $P/C$ the right side has no $C$ term. The lowering ideal is
nilpotent for the same bounded-weight reason as the raising ideal, so
iteration gives $P/C=0$. The proof needs full current stability of
$C$; the list of $28$ seeds itself need not be a current-stable basis.
The independent result is
\artifact{hi_cft_20260923_c8two_charged_span_verified.json}.

\subsubsection{How the neutral identity uses this particular spanning claim}

Because $P$ is generated by charged brackets, one may write
\[
 e_0a_3=\sum_jd_j(x_j)_0y_j+n,
 \qquad n\in N_3,
\]
where $x_j,y_j$ are charged fields, including their current
descendants. For the bottom input $p_3=\texttt{r3\_000910}$,
the zero-mode commutator gives
\[
 (e_0a_3)_0p_3
  =\sum_jd_j\bigl((x_j)_0(y_j)_0-(y_j)_0(x_j)_0\bigr)p_3
       +n_0p_3\in N_5.
\]
Here the inner products are in $W_4$ by Gate~4. The outer charged
products preserve membership modulo $N_5$ on its normal and triplet
parts by the charged conclusion and the normal-space identities.
A second application of Jacobi to $e,a_3,p_3$ then forces the neutral
bottom identity, since the coefficient of $P4_-$ in $(a_3)_0p_3$
is nonzero. Current covariance propagates it along the triplet, and
$f$ follows from the total ambient stress minus $e$ and the affine
Sugawara stresses.

This makes the dependencies of the saved Gate~5 flag explicit. The
record \artifact{hi_cft_20260923_c8two_jacobi_checked.json}
already has the exact small scalar data, but marks Gate~4 authentication,
charged-family derivation, neutral-family derivation and all-channel
completion false. Its status is conditional prerequisites. In contrast,
\artifact{hi_cft_20260923_c8two_gate5_verified.json}
authenticates Gate~4 and the charged-span certificate and records all
those flags true. The older conditional record is not evidence of a
contradiction: its prerequisites had not yet been supplied when it was
written.

\subsection{Why the unsuccessful direct routes remain in the archive}
\label{expfinite:historical-routes}

Several earlier files describe direct weight-five normal-product
searches, including the \texttt{n5\_direct}, \texttt{n5\_primitive},
\texttt{n5\_axisstop}, and \texttt{n5\_priority} families. Their
mathematical role was to search for positive containment identities
using a selected subset of actual normal words. They were useful
exploratory and engineering work. The completed Gate~5 conclusion uses
the Jacobi reduction above; it does not silently assume that all of
those direct searches completed.

The distinctions are especially relevant at weight four. The direct
pipeline retained $6007$ selected raw normal words, then adjoined the
actual triplet. Their characteristic-zero span had not been separately
shown to exhaust the full formal $N_4$. A prime rank over these columns
can establish a positive lower bound or a successful exact relation.
It cannot establish failure of containment in the full $N_4$, nor
prove a full characteristic-zero normal-space dimension. Consequently
``rank $6007$ plus three'' is not, by itself, a proof that
$\dim U_4=6010$.

The final strong-generation argument identifies $U_4=W_4$ by weight
bookkeeping, but the coefficient $6010$ in the full vacuum character
has a later character-theoretic dependency. This distinction is retained
even though the various numerical and modular counts are consistent
with $6010$.

The current-projector optimization called ``axis stop'' has a similarly
precise meaning. If a raising action is exactly zero on the current
vector, every remaining same-axis $FE$ factor acts as the identity.
The other commuting axis must still be processed. This is an algebraic
identity optimization and does not change projector ranges or permit
small numerical residuals to be treated as zero.

Finally, an apparently attractive shortcut through the single weak
generator $u$ was rejected in the source audits. Trying to show that
a Griess product preserves a not-yet-certified normal space introduces
\[
 (a_0x)_0p=a_0(x_0p)-x_0(a_0p).
\]
Even after Gate~4, $a_0p$ can contain a triplet component, so the second
term asks for exactly an unknown $x_0P_4$ identity. The eventual Jacobi
argument avoids this circle by choosing two concrete bottom witnesses,
using weight exclusion there, and only then propagating proved zero
identities by covariance. Weight exclusion is not asserted anew after
raising the bottom input.

\subsection{Finite membership of the structural weight-three fields}
\label{expfinite:structural-membership}

The Potts and parafermion primaries used later in the structural argument
also have finite generation obligations. It is not enough to recognize
their ambient OPEs: they must belong to the actual $U$. The portable
membership witness uses the same $120$ weight-two source fields and
$58$ accepted Cartan-zero weight-three words. Its source-word inventory
is
\[
34\text{ current normal products},\qquad
6\text{ translations},\qquad18\text{ zero-products}.
\]
The two target expansions have the following sizes:
\begin{center}\small
\begin{tabular}{@{}lrrrr@{}}
\toprule
Target&Ambient terms&Selected words&Rank at discovery&Pairs tried\\\midrule
Potts primary&$92$&$33$&$41$&$1$\\
Parafermion primary&$342$&$42$&$58$&$154$\\\bottomrule
\end{tabular}
\end{center}
Here ``pairs tried'' is the recorded cumulative number of Cartan-zero
weight-two pairs at discovery. The membership exporter has $43462$
persistent sparse terms in all of its source and target material.
The independent coordinate verifier decodes $21772$ basis rows,
checks every stored grade, confirms Cartan weight zero for all
weight-three sources and targets, and checks both exact coordinate sums.
These differing counts reflect different retained and decoded objects;
they are not competing claims for a vector-space dimension.

The complete expressions are in
\artifact{c8_reviewability_20260927/membership_witnesses.json}
and are printed in full, with their source-word conventions, in
Section~\ref{expfinite:complete-membership}. Its separate check record is
\artifact{c8_reviewability_20260927/membership_coordinate_check.json}.
That checker establishes saved coordinate sums and labels. It does not
recompute the source modes or prove the weight-two generation step;
those premises are supplied by the preceding finite construction.
The former rank-$41$ Potts membership assertion lacked a saved expression;
the $33$-word export is the additional concrete witness that repairs
that reviewability limitation.

\subsection{A compact provenance map for checking claims in this section}
\label{expfinite:provenance}

The next list orders the source records by their mathematical role,
not by filename or execution time. Every entry is a local artifact;
no claim of a new research replay is made by this transcription.
\begin{longtable}{@{}p{.34\linewidth}p{.60\linewidth}@{}}
\toprule Claim or object&Principal saved source\\\midrule\endfirsthead
\toprule Claim or object&Principal saved source\\\midrule\endhead
\bottomrule\endfoot
Ambient root enumeration and numerical primary kernels&
\artifact{hi_cft_20260923_c8_lattice_control.json}\\
Historical Virasoro-axis and primary-eigenvalue search&
\artifact{hi_cft_20260923_c8_ope_control.json}\\
All nine exact field-coordinate tables&
\artifact{hi_cft_20260923_c8_exact_fields.json}\\
Exact low-mode closure and transported PCT&
\artifact{hi_cft_20260923_c8_exact_unitary_certificate.json}\\
Short-generator membership graph&
\artifact{c8_reviewability_20260927/minimal_generation_witness.json}\\
Independent graph/coordinate verification&
\artifact{c8_reviewability_20260927/minimal_generation_checked.json}\\
Exact grade-three dimension and charge census&
\artifact{hi_cft_20260923_c8_grade3_exact_v2.json}\\
All lower-mode identities&
\artifact{hi_cft_20260923_next_structure_lower_closure.json}\\
Actual triplet, current action and Virasoro controls&
\artifact{hi_cft_20260923_next_structure_p4_controls.json}\\
Exact quotient raising maps and $130$ requests&
\artifact{hi_cft_20260923_c8two_structure_saturation.json}\\
Completed $118$ direct additional relations&
\artifact{hi_cft_20260923_c8two_targeted_h4_trimmed_v2/summary.json}\\
Completed twelve raised-input relations&
\artifact{hi_cft_20260923_c8two_structure_h4_raised_exact/summary.json}\\
Twenty-four raising premises for the singlet targets&
\artifact{hi_cft_20260923_c8two_h4_singlet_premise_verified.json}\\
Full projected exact residuals&
\artifact{hi_cft_20260923_c8two_h4_singlet_exact.json}\\
Portable residual-only independent replay&
\artifact{c8_reviewability_20260927/gate4_packet_checked.json}\\
Final Gate~4 recovery and saturation&
\artifact{hi_cft_20260923_c8two_gate4_verified.json}\\
Charged-span seeds and independent replay&
\artifact{hi_cft_20260923_c8two_structure_charged_span.json};
\artifact{hi_cft_20260923_c8two_charged_span_verified.json}\\
Final Gate~5 Jacobi prerequisites and conclusion&
\artifact{hi_cft_20260923_c8two_gate5_verified.json}\\
Potts/parafermion generated-field membership&
\artifact{c8_reviewability_20260927/membership_witnesses.json}\\
Repaired finite-input authentication ledger&
\artifact{c8_revision_20260927/finite_input_ledger.json}\\
\end{longtable}

The original proof-bundle ledger authenticated $640$ files but omitted
the independent short-generator record and its producer. That omission
mattered because the grade-four initializer imports saved highest
fields directly, whereas the short-generator record proves their
membership in $\langle K,\omega_B,u\rangle$. The independent
exact grade-three dimension record was also absent from that old ledger.
The $27$ September supplementary authentication bound these records,
the relevant fresh replay outputs, and their sources, reaching $650$
authenticated source/output files. Source authentication prevents
unnoticed changes of input; the equality and rank witnesses supply the
mathematical content.

For an efficient manual review, first check one full charged field from
its coordinate table and the elementary mode formula; then check the
short-generation coordinate expressions for that field; then inspect
its Gate~4 row and the common saturation argument. The remaining
Gate~5 work is the small bottom-witness/Jacobi calculation together with
the independent charged-span argument. This order keeps ambient
identities, actual generation, finite closure, and all-weight deductions
separate while making their connection explicit.


\subsection{Complete membership expressions for both structural primaries}
\label{expfinite:complete-membership}

Here are all $33$ Potts and all $42$ parafermion membership coefficients,
so neither expansion requires consulting the TSV file. The source is
\artifact{c8_reviewability_20260927/membership_witnesses.coefficients.tsv},
whose coefficients are also stored in the portable membership JSON.
The parafermion target is the \emph{unnormalized} saved primary $w$;
the normalized field $a=\sigma w$ used in the nonlinear OPE
presentation is obtained by multiplying the complete expression by
$\sigma$. The Potts target is the saved $w_P$.

The $S_2$ basis used here differs from the short-generation normalized
basis $r_i$. Denote it by $b_0,\ldots,b_{119}$. Its first $23$ vectors
are the \emph{raw} vectors $v_0,\ldots,v_{22}$ of
Section~\ref{expfinite:generation-graph}; the saved coordinates agree
entry by entry. Next, $b_{23}=e$. The remaining $96$ are the charged
multiplets, in the following completely specified order:
\[
 b_{q_i+m_i\ell+j}=(F_L)_0^\ell(F_R)_0^ja_i,
 \qquad0\le\ell<n_i,\quad0\le j<m_i.
\]
There are no pivot normalizations or divided powers in this formula.
\begin{center}\small
\begin{tabular}{@{}rrrrl@{}}
\toprule
$i$&$q_i$&$n_i$&$m_i$&Range of $b$ indices\\\midrule
0&24&1&5&24--28\\
1&29&1&9&29--37\\
2&38&1&9&38--46\\
3&47&1&13&47--59\\
4&60&2&2&60--63\\
5&64&2&6&64--75\\
6&76&2&8&76--91\\
7&92&2&14&92--119\\\bottomrule
\end{tabular}
\end{center}
Thus $b_{23}=e$, whereas $r_{23}=v_{23}=u$ in the short-generation
graph. This single example makes the index mismatch visible. The
checkpoint labels beginning \texttt{r2\_} provide yet a third
basis convention and are not used in the tables below.

Write $z_j$ for the indicated unreduced weight-three source word.
Each table asserts the actual lattice-vector equality
$\text{target}=\sum_j c_jz_j$, with precisely the listed nonzero
coefficients. A current factor such as $(H_L)_{-1}b_0$ is a normal
product; $Tb_i$ is an ambient translation; $(b_i)_0b_k$ is a
zero-product. The source ID $j$ is retained solely to allow a direct
comparison with the exported certificate.

\subsubsection{Potts primary}
\begingroup\footnotesize\renewcommand{\arraystretch}{1.55}
\begin{longtable}{@{}rlr@{}}
\toprule Source $j$&Actual source word $z_j$&Coefficient $c_j$\\\midrule\endfirsthead
\toprule Source $j$&Actual source word $z_j$&Coefficient $c_j$\\\midrule\endhead
\bottomrule\endfoot
0 & $(H_L)_{-1}b_{0}$ & $\frac{1053}{70}$\\
1 & $(H_L)_{-1}b_{1}$ & $\frac{-729}{280}$\\
2 & $(H_L)_{-1}b_{4}$ & $\frac{1647}{280}$\\
3 & $(H_L)_{-1}b_{13}$ & $\frac{-190971}{27265}$\\
4 & $(H_L)_{-1}b_{14}$ & $\frac{-204741}{218120}$\\
5 & $(H_L)_{-1}b_{19}$ & $\frac{27729}{54530}$\\
6 & $(H_L)_{-1}b_{22}$ & $\frac{-1053}{35}$\\
7 & $(H_L)_{-1}b_{23}$ & $\frac{2511}{140}$\\
8 & $(H_L)_{-1}b_{26}$ & $\frac{117}{152}$\\
9 & $(H_L)_{-1}b_{33}$ & $\frac{-27s-54t}{35200}$\\
10 & $(H_L)_{-1}b_{42}$ & $\frac{9s-27t}{44800}$\\
11 & $(H_L)_{-1}b_{53}$ & $\frac{-2+t}{985600}$\\
12 & $(E_L)_{-1}b_{3}$ & $\frac{81}{2}$\\
13 & $(E_L)_{-1}b_{10}$ & $\frac{-27}{2}$\\
14 & $(H_R)_{-1}b_{13}$ & $\frac{60903}{21812}$\\
15 & $(H_R)_{-1}b_{14}$ & $\frac{-13717}{218120}$\\
16 & $(H_R)_{-1}b_{19}$ & $\frac{-118071}{27265}$\\
17 & $(H_R)_{-1}b_{22}$ & $\frac{-351}{35}$\\
18 & $(H_R)_{-1}b_{23}$ & $\frac{837}{140}$\\
19 & $(H_R)_{-1}b_{26}$ & $\frac{39}{152}$\\
20 & $(H_R)_{-1}b_{33}$ & $\frac{-9s-18t}{35200}$\\
21 & $(H_R)_{-1}b_{42}$ & $\frac{3s-9t}{44800}$\\
22 & $(H_R)_{-1}b_{53}$ & $\frac{-2+t}{2956800}$\\
23 & $(E_R)_{-1}b_{16}$ & $\frac{-27729}{109060}$\\
24 & $(E_R)_{-1}b_{20}$ & $-9$\\
29 & $(F_R)_{-1}b_{15}$ & $\frac{518499}{109060}$\\
34 & $Tb_{22}$ & $\frac{-1053}{35}$\\
35 & $Tb_{23}$ & $\frac{2511}{140}$\\
36 & $Tb_{26}$ & $\frac{117}{152}$\\
37 & $Tb_{33}$ & $\frac{-27s-54t}{35200}$\\
38 & $Tb_{42}$ & $\frac{9s-27t}{44800}$\\
39 & $Tb_{53}$ & $\frac{-2+t}{985600}$\\
40 & $(b_{92})_0b_{119}$ & $\frac{-1}{3985293312000}$\\
\end{longtable}
\endgroup
\subsubsection{Unnormalized parafermion primary}
\begingroup\footnotesize\renewcommand{\arraystretch}{1.55}
\begin{longtable}{@{}rlr@{}}
\toprule Source $j$&Actual source word $z_j$&Coefficient $c_j$\\\midrule\endfirsthead
\toprule Source $j$&Actual source word $z_j$&Coefficient $c_j$\\\midrule\endhead
\bottomrule\endfoot
0 & $(H_L)_{-1}b_{0}$ & $\frac{-4901s-14703t}{280}$\\
1 & $(H_L)_{-1}b_{1}$ & $\frac{1521s+4563t}{1120}$\\
2 & $(H_L)_{-1}b_{4}$ & $\frac{-13637s-40911t}{3360}$\\
3 & $(H_L)_{-1}b_{13}$ & $\frac{1195337s+3586011t}{327180}$\\
4 & $(H_L)_{-1}b_{14}$ & $\frac{1281527s+3844581t}{2617440}$\\
5 & $(H_L)_{-1}b_{19}$ & $\frac{-173563s-520689t}{654360}$\\
6 & $(H_L)_{-1}b_{22}$ & $\frac{2197s+6591t}{140}$\\
7 & $(H_L)_{-1}b_{23}$ & $\frac{-5239s-15717t}{560}$\\
8 & $(H_L)_{-1}b_{26}$ & $\frac{-2197s-6591t}{5472}$\\
9 & $(H_L)_{-1}b_{33}$ & $\frac{-169+169st}{84480}$\\
10 & $(H_L)_{-1}b_{42}$ & $\frac{-169}{40320}$\\
11 & $(H_L)_{-1}b_{53}$ & $\frac{1521+338s+1014t-169st}{319334400}$\\
12 & $(E_L)_{-1}b_{3}$ & $\frac{-169s-507t}{8}$\\
13 & $(E_L)_{-1}b_{10}$ & $\frac{169s+507t}{24}$\\
14 & $(H_R)_{-1}b_{13}$ & $\frac{-729383s-2188149t}{261744}$\\
15 & $(H_R)_{-1}b_{14}$ & $\frac{788891s+2366673t}{34026720}$\\
16 & $(H_R)_{-1}b_{19}$ & $\frac{398907s+1196721t}{109060}$\\
17 & $(H_R)_{-1}b_{22}$ & $\frac{1287s+3861t}{140}$\\
18 & $(H_R)_{-1}b_{23}$ & $\frac{403s+1209t}{560}$\\
19 & $(H_R)_{-1}b_{26}$ & $\frac{-1079s-3237t}{5472}$\\
20 & $(H_R)_{-1}b_{33}$ & $\frac{-3+3st}{2560}$\\
21 & $(H_R)_{-1}b_{42}$ & $\frac{29}{40320}$\\
22 & $(H_R)_{-1}b_{53}$ & $\frac{9+2s+6t-st}{3225600}$\\
23 & $(E_R)_{-1}b_{16}$ & $\frac{318929s+956787t}{1308720}$\\
24 & $(E_R)_{-1}b_{20}$ & $\frac{91s+273t}{12}$\\
29 & $(F_R)_{-1}b_{15}$ & $\frac{-5281159s-15843477t}{1308720}$\\
34 & $Tb_{22}$ & $\frac{4901s+14703t}{140}$\\
35 & $Tb_{23}$ & $\frac{5239s+15717t}{560}$\\
36 & $Tb_{26}$ & $\frac{-2866747s-8600241t}{5296896}$\\
37 & $Tb_{33}$ & $\frac{-111163+111163st}{40888320}$\\
38 & $Tb_{42}$ & $\frac{-377}{19514880}$\\
39 & $Tb_{53}$ & $\frac{1453959+323102s+969306t-161551st}{309115699200}$\\
40 & $(b_{92})_0b_{119}$ & $\frac{13s+39t}{99325771776000}$\\
49 & $(b_{49})_0b_{37}$ & $\frac{-26+39s+13t+26st}{37468569600}$\\
50 & $(b_{49})_0b_{46}$ & $\frac{26-13t}{2554675200}$\\
51 & $(b_{49})_0b_{57}$ & $\frac{-468-13s-39t+52st}{303495413760000}$\\
52 & $(b_{57})_0b_{29}$ & $\frac{26-39s-13t-26st}{1686085632000}$\\
53 & $(b_{57})_0b_{38}$ & $\frac{-26+13t}{114960384000}$\\
54 & $(b_{29})_0b_{37}$ & $\frac{91s-143t}{1040793600}$\\
55 & $(b_{29})_0b_{46}$ & $\frac{-13s-26t}{70963200}$\\
56 & $(b_{37})_0b_{38}$ & $\frac{13s+26t}{70963200}$\\
57 & $(b_{38})_0b_{46}$ & $\frac{-13s+39t}{38707200}$\\
\end{longtable}
\endgroup
Each word on the right is in $U_3$ because every $b_i$ is in $U_2$
and the currents, translation and vertex modes preserve $U$. Thus,
once the mode-word provenance and exact coordinate sum are accepted,
these two tables prove membership of the actual saved primary fields.
The OPE identifications of those fields remain separate structural
calculations; a membership expression alone does not identify a Potts
or parafermion algebra.


\subsection{Historical discovery spectra, retained as numerical diagnostics}
\label{expfinite:discovery-spectra}

For completeness, the following two tables retain the exploratory
spectral information in
\artifact{hi_cft_20260923_c8_ope_control.json}.
They are rounded decimal transcriptions of that floating run and carry
no exact status. Their usefulness is diagnostic: a reviewer can see the
candidate eigenvalues available before the exact field selection, and
can distinguish a repeated numerical eigenvalue from a selected vector.

The commutant Griess calculation reported the following sorted list of
axis central charges. The saved record calls this list
\texttt{all\_axis\_charges}; this is the output of that search, rather
than a new assertion here that every possible Virasoro axis has been
classified. The selected stress has exact charge $2400/1271$ and its
commuting complement charge $70/31$ in the later exact certificate.

\begin{longtable}{@{}rr@{}}
\toprule Sorted index&Numerical axis charge\\\midrule\endfirsthead
\toprule Sorted index&Numerical axis charge\\\midrule\endhead
\bottomrule\endfoot
0 & $0$\\
1 & $0.799999999999999$\\
2 & $1.36216216216216$\\
3 & $1.80343063282447$\\
4 & $1.80343063282447$\\
5 & $1.88827694728558$\\
6 & $1.8882769472856$\\
7 & $1.98417930125248$\\
8 & $2.16216216216196$\\
9 & $2.25806451612902$\\
10 & $2.25806451612904$\\
11 & $2.34291083059015$\\
12 & $2.34291083059016$\\
13 & $2.78417930125247$\\
14 & $3.34634146341464$\\
15 & $4.14634146341463$\\
\end{longtable}
The next table gives all eigenvalues of the selected numerical
$e_1$ on each relevant ambient affine-highest space. The final column
is $\|e_2v\|$ for the corresponding normalized numerical eigenvector.
It was computed in the original orthonormal oscillator convention, so
it should not be compared directly with an untransported norm of the
exact-gauge coordinate table. An eigenvector with a visibly nonzero
final entry is not primary for this Virasoro stress, even though it
belongs to the affine-highest space. Tiny entries are floating residuals
and do not themselves prove exact vanishing.

For repeated numerical eigenvalues the particular eigenbasis returned
by a diagonalizer need not equal the later exact chosen basis. In
particular, multiplicity and phase information cannot be recovered
from this list alone. The exact charged fields, their $e_1$
eigenvalues, and their positive-mode checks supply the subsequent
selection.

\begingroup\small
\begin{longtable}{@{}crrr@{}}
\toprule Current type&Eigenvector index&Numerical $e_1$ eigenvalue&$\|e_2v\|$\\\midrule\endfirsthead
\toprule Current type&Eigenvector index&Numerical $e_1$ eigenvalue&$\|e_2v\|$\\\midrule\endhead
\bottomrule\endfoot
$(1,5)$ & 0 & $0.0524850686944$ & $9.08421\times10^{-16}$\\
$(1,5)$ & 1 & $0.122084992219$ & $3.92597\times10^{-15}$\\
$(1,5)$ & 2 & $0.26983331241$ & $6.95716\times10^{-16}$\\
$(1,5)$ & 3 & $0.627852084972$ & $1.44018\times10^{-14}$\\
$(1,5)$ & 4 & $0.627852084972$ & $8.67494\times10^{-16}$\\
$(1,5)$ & 5 & $0.714264339016$ & $1.60186\times10^{-14}$\\
$(1,5)$ & 6 & $0.966153927523$ & $2.24913\times10^{-15}$\\
$(1,5)$ & 7 & $1.12208499222$ & $0.494135593171$\\
$(1,5)$ & 8 & $1.46124613193$ & $2.33481\times10^{-16}$\\
$(1,5)$ & 9 & $1.71426433902$ & $1.19521072537$\\
$(1,5)$ & 10 & $1.7766387585$ & $9.22827\times10^{-15}$\\
$(1,9)$ & 0 & $0.0881552109736$ & $2.23261\times10^{-16}$\\
$(1,9)$ & 1 & $0.157356412274$ & $6.10418\times10^{-16}$\\
$(1,9)$ & 2 & $0.157356412274$ & $3.21692\times10^{-15}$\\
$(1,9)$ & 3 & $0.318646734854$ & $4.48097\times10^{-15}$\\
$(1,9)$ & 4 & $0.580108555496$ & $2.60578\times10^{-16}$\\
$(1,9)$ & 5 & $0.738001573564$ & $6.55504\times10^{-15}$\\
$(1,9)$ & 6 & $0.775315891617$ & $6.97350\times10^{-17}$\\
$(1,9)$ & 7 & $1.18961447679$ & $2.41474\times10^{-15}$\\
$(1,9)$ & 8 & $1.18961447679$ & $4.12746\times10^{-15}$\\
$(1,9)$ & 9 & $1.31864673485$ & $0.798306626372$\\
$(1,9)$ & 10 & $1.44391050714$ & $3.10381\times10^{-15}$\\
$(1,13)$ & 0 & $0.0930995402395$ & $0$\\
$(1,13)$ & 1 & $0.330448465775$ & $0$\\
$(1,13)$ & 2 & $0.330448465775$ & $0$\\
$(1,13)$ & 3 & $0.556254917388$ & $0$\\
$(1,13)$ & 4 & $0.82586190744$ & $0$\\
$(2,2)$ & 0 & $0.151062155783$ & $1.97606\times10^{-16}$\\
$(2,2)$ & 1 & $0.151062155783$ & $3.36752\times10^{-16}$\\
$(2,2)$ & 2 & $1.05428796223$ & $2.53151\times10^{-16}$\\
$(2,2)$ & 3 & $1.53815892998$ & $1.13555\times10^{-16}$\\
$(2,6)$ & 0 & $0.0701906421897$ & $2.02179\times10^{-15}$\\
$(2,6)$ & 1 & $0.166719040056$ & $6.82318\times10^{-16}$\\
$(2,6)$ & 2 & $0.310778914241$ & $1.67610\times10^{-15}$\\
$(2,6)$ & 3 & $0.64814137984$ & $9.57819\times10^{-15}$\\
$(2,6)$ & 4 & $0.874193627135$ & $6.99816\times10^{-16}$\\
$(2,6)$ & 5 & $1.31077891424$ & $0.788389388869$\\
$(2,6)$ & 6 & $1.43981117231$ & $7.96599\times10^{-16}$\\
$(2,6)$ & 7 & $1.43981117231$ & $6.90936\times10^{-15}$\\
$(2,8)$ & 0 & $0.052056041325$ & $2.08167\times10^{-17}$\\
$(2,8)$ & 1 & $0.164636264896$ & $6.59195\times10^{-17}$\\
$(2,8)$ & 2 & $0.591660110149$ & $1.73472\times10^{-16}$\\
$(2,8)$ & 3 & $0.591660110149$ & $6.93889\times10^{-17}$\\
$(2,8)$ & 4 & $0.591660110149$ & $2.94903\times10^{-16}$\\
$(2,8)$ & 5 & $0.986425890887$ & $2.77556\times10^{-17}$\\
$(2,8)$ & 6 & $1.22803837252$ & $6.24500\times10^{-17}$\\
$(2,14)$ & 0 & $0.141620771046$ & $0$\\
$(2,14)$ & 1 & $0.141620771046$ & $0$\\
$(2,14)$ & 2 & $0.496459480724$ & $0$\\
\end{longtable}\endgroup
The largest leakage of $e_1$ out of these affine-highest spaces
is on the order of $10^{-15}$ in this saved run. The two conspicuously
nonzero $e_2$ norms in the $(1,5)$ space, and the nonzero norms in the
$(1,9)$ and $(2,6)$ spaces, explain why diagonalizing $e_1$ alone
would have been insufficient. The exact primary checks test both the
eigenvalue equation and the required positive-mode annihilation.

